\documentclass{amsart}
\usepackage{amsmath, amssymb, amsthm, hyperref, tikz}

\title{The Kreuzer--Skarke Classification Programme \\
       for Reflexive Polytopes}
\author{Personalized Notes (companion to \texttt{refpoly-5d})}
\date{\today}

\newtheorem{theorem}{Theorem}[section]
\newtheorem{lemma}[theorem]{Lemma}
\newtheorem{proposition}[theorem]{Proposition}
\newtheorem{corollary}[theorem]{Corollary}
\newtheorem{definition}[theorem]{Definition}
\newtheorem{example}[theorem]{Example}
\newtheorem{remark}[theorem]{Remark}

\begin{document}

\begin{abstract}
These notes are a standalone companion to \texttt{theory\_notes.tex}. They
develop, from first principles, the structural backbone of the
Kreuzer--Skarke (KS) classification of reflexive polytopes: why the count
is finite in each dimension, why every reflexive polytope sits inside a
\emph{maximal} one, how maximal polytopes are polar-dual to \emph{minimal}
ones, how minimal polytopes decompose as simplices or products of
simplices, and how these are bounded and constructed from \emph{weight
systems}. Throughout we keep the eventual computational target in view:
the complete $d=4$ enumeration of $473\,800\,776$ reflexive polytopes, and
the open frontier in $d=5$ addressed by \texttt{refpoly-5d}.
We use the conventions of the companion notes: dual lattices
$M,N \cong \mathbb{Z}^d$, real extensions $M_{\mathbb{R}},N_{\mathbb{R}}$,
pairing $\langle\,\cdot\,,\,\cdot\,\rangle$, and polar dual $\Delta^*$
(Definition in the companion document, Batyrev 1994). Recall a lattice
polytope $\Delta \subset M_{\mathbb{R}}$ with $0$ in its interior is
\emph{reflexive} iff $\Delta^*$ is also a lattice polytope, equivalently iff
$\Delta$ has exactly one interior lattice point.
\end{abstract}

\maketitle
\tableofcontents

%==============================================================
\section{Overview of the programme and its computational shape}
%==============================================================

The Kreuzer--Skarke programme is the answer to a deceptively simple
question: \emph{list all reflexive polytopes in dimension $d$, up to
$GL(d,\mathbb{Z})$ equivalence}. The completed tally is
\begin{equation}\label{eq:counts}
  d=1:\ 1, \qquad d=2:\ 16, \qquad d=3:\ 4319,
  \qquad d=4:\ 473\,800\,776,
\end{equation}
with $d=5$ presently unclassified (the target of \texttt{refpoly-5d}).
The explosive growth in~\eqref{eq:counts} is precisely why a naive
enumeration is hopeless and why the programme is organized around the
following chain of reductions, each a section below:

\begin{enumerate}
  \item[(F)] \textbf{Finiteness.} There are only finitely many reflexive
        polytopes in each dimension (Section~\ref{sec:finite}). This makes
        the problem well-posed and gives crude a~priori size bounds.
  \item[(S)] \textbf{Subpolytope reduction.} Every reflexive polytope is a
        subpolytope of a \emph{maximal} reflexive polytope
        (Section~\ref{sec:sub}). So it suffices to find all maximal ones
        and enumerate their reflexive subpolytopes.
  \item[(D)] \textbf{Max/min duality.} Maximal reflexive polytopes are
        exactly the polar duals of \emph{minimal} reflexive polytopes
        (Section~\ref{sec:dual}). So it suffices to find all minimal ones.
  \item[(M)] \textbf{Decomposition of minimal polytopes.} Minimal reflexive
        polytopes are simplices, or convex hulls built from
        lower-dimensional simplices ("products of simplices")
        (Section~\ref{sec:min}).
  \item[(W)] \textbf{Weight systems.} The minimal building blocks are
        enumerated from \emph{weight systems} $q=(q_0,\dots,q_d)$ with the
        interior-point property; products of simplices come from
        \emph{combined} weight systems (Section~\ref{sec:weights}).
\end{enumerate}

\noindent The actual KS computation runs this chain \emph{in reverse}:
enumerate weight systems $\to$ build minimal polytopes $\to$ dualize to
get maximal polytopes $\to$ enumerate reflexive subpolytopes $\to$
normal-form and deduplicate. The \texttt{refpoly-5d} pipeline implements
the bracketing steps (weight-system input, $GL(d,\mathbb{Z})$ normal forms
via PALP, and deduplication) for $d=5$.

%==============================================================
\section{Finiteness}\label{sec:finite}
%==============================================================

\subsection{The statement}

\begin{theorem}[Finiteness; Lagarias--Ziegler 1991]\label{thm:finite}
For each fixed $d$, the number of reflexive polytopes in $\mathbb{R}^d$,
up to the action of $GL(d,\mathbb{Z})$ (and lattice translation), is finite.
\end{theorem}

This is the indispensable foundation: without it, "classify all reflexive
polytopes in dimension $d$" would not even be a finite problem. It rests on
two pillars --- a \emph{volume bound} for lattice polytopes with interior
points, and a \emph{boundedness} (up to $GL(d,\mathbb{Z})$) of lattice
polytopes of bounded volume.

\subsection{Volume bounds: Hensley and Lagarias--Ziegler}

The first pillar controls how big a reflexive polytope can be.

\begin{theorem}[Hensley 1983; Lagarias--Ziegler 1991]\label{thm:hensley}
There is a constant $c(d,k)$ such that any $d$-dimensional lattice polytope
$P \subset \mathbb{R}^d$ with exactly $k \geq 1$ interior lattice points has
normalized volume
\begin{equation}
  \operatorname{Vol}(P) \;\leq\; c(d,k).
\end{equation}
In particular a reflexive polytope ($k=1$ interior point) has volume bounded
by a constant depending only on $d$.
\end{theorem}

\begin{remark}[Why an interior point forces boundedness]
The intuition is a "no room to grow" principle. If a lattice polytope $P$
has an interior lattice point $p$ but is very long and thin (large volume,
few interior points), then by a pigeonhole/convexity argument one can
exhibit \emph{additional} interior lattice points: a sufficiently elongated
lattice polytope around $p$ must swallow a lattice point in its relative
interior in some direction. Hensley's theorem quantifies this: fixing the
number of interior points $k$ caps the volume. The original Hensley bound was
doubly exponential; Lagarias--Ziegler sharpened the argument and, crucially
for us, drew the classification consequence. Explicit values are large
(e.g.\ the maximal reflexive simplex volumes grow like the Sylvester
sequence; in $d$ dimensions the largest reflexive simplex is associated with
weights $(1,1,2,3,7,42,\dots)$-type "Sylvester" data), but \emph{finite}.
\end{remark}

\subsection{From volume bound to finiteness}

The second pillar converts "bounded volume" into "finitely many up to
$GL(d,\mathbb{Z})$".

\begin{lemma}[Bounded volume $\Rightarrow$ finiteness up to $GL(d,\mathbb{Z})$]
\label{lem:boundedfinite}
For fixed $d$ and fixed bound $V$, there are only finitely many lattice
polytopes $P \subset \mathbb{R}^d$ with $\operatorname{Vol}(P) \leq V$ and a
lattice point in the interior, up to $GL(d,\mathbb{Z})$ and translation.
\end{lemma}

\begin{proof}[Mechanism]
Translate the interior lattice point to the origin. A theorem in the
geometry of numbers (a quantitative form of "a lattice polytope of bounded
volume can be put into a box of bounded size") shows $P$ is
$GL(d,\mathbb{Z})$-equivalent to a polytope contained in a fixed finite box
$[-R,R]^d$ with $R=R(d,V)$. Lattice polytopes inside a fixed finite box are
finite in number (their vertex sets are subsets of $[-R,R]^d \cap
\mathbb{Z}^d$), so finitely many $GL(d,\mathbb{Z})$-classes remain.
\end{proof}

\begin{proof}[Proof of Theorem~\ref{thm:finite}]
A reflexive polytope has exactly one interior lattice point, so by
Theorem~\ref{thm:hensley} its volume is at most $c(d,1)$. Apply
Lemma~\ref{lem:boundedfinite} with $V=c(d,1)$.
\end{proof}

\begin{remark}[The role of the dual polytope]
Reflexivity is self-dual: $\Delta$ is reflexive iff $\Delta^*$ is, and
$(\Delta^*)^* = \Delta$. Finiteness therefore comes in dual pairs, and
\emph{both} $\Delta$ and $\Delta^*$ are simultaneously controlled by
Theorem~\ref{thm:hensley}. This is more than cosmetic: the KS strategy
exploits the fact that bounding one member of the pair bounds the other,
and that the "large" member ($\Delta^*$ when $\Delta$ is "small") is the one
with rich lattice-point structure. The volume bound on $\Delta^*$ caps the
number of lattice points $\Delta^* \cap N$, which (via Batyrev) caps
$h^{1,1}$ of the associated Calabi--Yau and the number of monomials in the
anticanonical hypersurface.
\end{remark}

\begin{remark}[Why finiteness is not enough]
The bounds $c(d,1)$ are astronomically large, so Theorem~\ref{thm:finite}
gives no practical algorithm. The counts in~\eqref{eq:counts} were obtained
not by brute force inside a box, but by the structural reductions
(S)--(W) below, which is the genuine content of the KS programme.
\end{remark}

%==============================================================
\section{Subpolytope structure: maximal reflexive polytopes}\label{sec:sub}
%==============================================================

\begin{definition}[Subpolytope; maximal reflexive polytope]\label{def:max}
For reflexive polytopes $P,\Delta$ in the \emph{same} lattice $M$, we say
$P$ is a \emph{(reflexive) subpolytope} of $\Delta$, written $P \subseteq
\Delta$, if every vertex of $P$ lies in $\Delta \cap M$ (equivalently
$P \subseteq \Delta$ as point sets, both having the single interior point
$0$). A reflexive polytope $\Delta$ is \emph{maximal} if it is not a proper
reflexive subpolytope of any other reflexive polytope --- i.e.\ one cannot
enlarge $\Delta$ to a strictly larger reflexive polytope on the same lattice.
\end{definition}

\begin{theorem}[Every reflexive polytope embeds in a maximal one]\label{thm:embed}
Every reflexive polytope $P \subset M_{\mathbb{R}}$ is a subpolytope of some
maximal reflexive polytope $\Delta \subset M_{\mathbb{R}}$.
\end{theorem}

\begin{proof}[Mechanism]
Order reflexive polytopes on the lattice $M$ by inclusion. By
Theorem~\ref{thm:finite} there are only finitely many, so any ascending
chain $P = P_0 \subsetneq P_1 \subsetneq \cdots$ terminates; a maximal
element $\Delta$ above $P$ exists. Concretely: while $P$ is not maximal there
is a lattice point $v$ with $\operatorname{conv}(P \cup \{v\})$ still
reflexive; adjoin it and repeat. Finiteness guarantees the process halts.
\end{proof}

\begin{remark}[Why this is the central reduction]
Theorem~\ref{thm:embed} reduces the whole classification to two tasks:
(i) enumerate the \emph{maximal} reflexive polytopes; (ii) for each maximal
$\Delta$, enumerate the reflexive subpolytopes $P \subseteq \Delta$. Step
(ii) is a finite, downward search inside the lattice points of a single
$\Delta$: pick subsets of $\Delta \cap M$ whose convex hull is reflexive.
The payoff is enormous because the maximal polytopes are comparatively
few. The dual statement is the actual KS workhorse: passing to duals,
$P \subseteq \Delta$ becomes $\Delta^* \subseteq P^*$, so subpolytopes of a
maximal $\Delta$ correspond to polytopes \emph{containing} the
corresponding minimal $\Delta^*$ --- which is how the next section reframes
everything.
\end{remark}

%==============================================================
\section{Duality of maximal and minimal polytopes}\label{sec:dual}
%==============================================================

Polar duality reverses inclusion:
\begin{equation}\label{eq:incl-rev}
  P \subseteq \Delta \quad\Longleftrightarrow\quad \Delta^* \subseteq P^*,
\end{equation}
because enlarging a polytope shrinks the set of linear functionals bounded
below by $-1$ on it. This single fact turns "maximal" into "minimal".

\begin{definition}[Minimal reflexive polytope]\label{def:min}
A reflexive polytope $\nabla \subset N_{\mathbb{R}}$ is \emph{minimal} if it
contains no proper reflexive subpolytope, equivalently if removing any
vertex destroys the interior-point / reflexivity property: $\nabla$ is a
smallest reflexive polytope still wrapping the origin.
\end{definition}

\begin{theorem}[Max/min duality]\label{thm:maxmin}
Polar duality is an inclusion-reversing bijection between reflexive
polytopes in $M$ and reflexive polytopes in $N$. Under it:
\begin{equation}
  \Delta \text{ maximal in } M
  \quad\Longleftrightarrow\quad
  \nabla := \Delta^* \text{ minimal in } N.
\end{equation}
Consequently, every reflexive polytope $P$ lies in a "family" determined by
a minimal polytope $\nabla$: namely $P^* \supseteq \nabla = \Delta^*$ where
$\Delta \supseteq P$ is a maximal hull of $P$ from
Theorem~\ref{thm:embed}.
\end{theorem}

\begin{proof}
Duality is a bijection on reflexive polytopes (involutivity, companion
notes). By~\eqref{eq:incl-rev} it reverses inclusion, hence sends maximal
elements to minimal elements and vice versa. The last clause is
Theorem~\ref{thm:embed} dualized through~\eqref{eq:incl-rev}.
\end{proof}

\begin{remark}[The strategic inversion]
Maximal polytopes are hard to find directly (they are "big": many lattice
points, complicated facets). Minimal polytopes are easy to find directly
(they are "small": few vertices, tightly constrained by the requirement
that $0$ be interior). Theorem~\ref{thm:maxmin} says we may enumerate the
\emph{easy} objects and dualize. This is the conceptual pivot of KS: the
classification of all reflexive polytopes is bootstrapped from the
classification of the minimal ones, which --- as the next two sections show
--- are governed by simple combinatorial-arithmetic data (weight systems).
\end{remark}

\begin{remark}[Picture in $d=2$]
The minimal reflexive polygons are the three triangles dual to the maximal
triangles/quadrilaterals. For instance the minimal $\nabla$ with vertices
$(1,0),(0,1),(-1,-1)$ (the $\mathbb{P}^2$ fan generators) is dual to the
maximal reflexive triangle with vertices $(-1,-1),(2,-1),(-1,2)$. All $16$
reflexive polygons arise as subpolytopes of a handful of maximal ones, dual
to a handful of minimal triangles/segments-products.
\end{remark}

%==============================================================
\section{Decomposition of minimal reflexive polytopes}\label{sec:min}
%==============================================================

We now describe the minimal polytopes $\nabla \subset N_{\mathbb{R}}$
concretely. The defining feature is that $0$ is an interior point obtained
"as economically as possible". This forces a rigid structure.

\begin{definition}[IP (interior point) property]\label{def:ip}
A finite set of lattice vectors $V = \{v_0,\dots,v_k\} \subset N$ has the
\emph{IP property} if $0 \in \operatorname{int}\,\operatorname{conv}(V)$,
i.e.\ the origin is an interior point of their convex hull. A minimal
reflexive polytope is the convex hull of an IP set that is minimal for
inclusion (no proper subset is IP).
\end{definition}

\begin{theorem}[Structure of minimal IP configurations]\label{thm:minstructure}
Let $\nabla = \operatorname{conv}(V) \subset N_{\mathbb{R}}$ be a minimal
reflexive polytope, $\dim \nabla = d$. Then the vertex set $V$ decomposes as
a disjoint union of "simplicial blocks"
\begin{equation}
  V = V_1 \sqcup V_2 \sqcup \cdots \sqcup V_r,
\end{equation}
where each block $V_j$ is an affinely independent set of $d_j+1$ vectors
spanning a $d_j$-dimensional sublattice, $\sum_j d_j = d$, and $0$ lies in
the relative interior of $\operatorname{conv}(V_j)$ within that block's span.
Equivalently:
\begin{itemize}
  \item if $r=1$, $\nabla$ is a \emph{simplex} (the simplicial case): $d+1$
        vertices in general position with $0$ interior;
  \item if $r>1$, $\nabla$ is a \emph{convex hull of a collection of
        lower-dimensional simplices} whose spans are complementary and meet
        only at $0$ --- a "product of simplices" structure.
\end{itemize}
\end{theorem}

\begin{remark}[Why minimality forces this]
For $0$ to be interior to $\operatorname{conv}(V)$ one needs a positive
combination $\sum_i \lambda_i v_i = 0$, $\lambda_i>0$, spanning all $d$
dimensions. Minimality (no proper IP subset) means $V$ admits no
\emph{smaller} such positive dependence on a subset spanning the same space.
Carath\'eodory-type reasoning then shows $V$ breaks into independent
positive circuits: each block $V_j$ is a single minimal positive dependence
(a simplex with $0$ in its interior), and distinct blocks live in
complementary subspaces. A single circuit gives a simplex; several stacked
in complementary directions give the product-of-simplices shape. This is
exactly the combinatorial heart of Kreuzer--Skarke (1998, 2000): the minimal
polytopes are classified by how the $d$ dimensions are partitioned into
simplicial circuits.
\end{remark}

\begin{example}[Simplex vs.\ product in low dimension]\label{ex:minexamples}
\leavevmode
\begin{itemize}
  \item \emph{Simplicial ($r=1$, $d=2$):} $V=\{(1,0),(0,1),(-1,-1)\}$,
        the $\mathbb{P}^2$ triangle. One circuit
        $1\cdot(1,0)+1\cdot(0,1)+1\cdot(-1,-1)=0$.
  \item \emph{Product ($r=2$, $d=2$):} $V=\{(\pm1,0),(0,\pm1)\}$, the
        "cross" of $\mathbb{P}^1\times\mathbb{P}^1$. Two
        one-dimensional circuits $(1,0)+(-1,0)=0$ and $(0,1)+(0,-1)=0$ in
        complementary lines; $\nabla$ is the square (a product of two
        $1$-simplices).
\end{itemize}
Dually, these are the minimal polytopes whose maximal duals seed the
$\mathbb{P}^2$ and $\mathbb{P}^1\times\mathbb{P}^1$ families of reflexive
polygons.
\end{example}

%==============================================================
\section{Weight systems}\label{sec:weights}
%==============================================================

The final layer makes Section~\ref{sec:min} arithmetic and therefore
enumerable. Each simplicial block is encoded by a \emph{weight system}.

\subsection{Single weight systems and the simplicial case}

\begin{definition}[Weight system]\label{def:weights}
A \emph{weight system} of length $d+1$ is a tuple of positive numbers
\begin{equation}
  q = (q_0, q_1, \dots, q_d), \qquad q_i \in \mathbb{Q}_{>0},
\end{equation}
normalized by $\sum_{i=0}^{d} q_i = 1$. Equivalently one uses coprime
positive \emph{integers} $(w_0,\dots,w_d)$ with $q_i = w_i / \sum_j w_j$;
the common denominator $\sum_j w_j$ is the \emph{degree} $d_{\mathrm{w}}$.
\end{definition}

A weight system specifies a circuit: a simplex with vertices
$v_0,\dots,v_d$ obeying the single linear relation
\begin{equation}\label{eq:circuit}
  \sum_{i=0}^{d} w_i\, v_i = 0,
\end{equation}
which is exactly the condition that $0 \in N_{\mathbb{R}}$ lie in the
relative interior with barycentric-type coordinates proportional to the
weights. Geometrically $q$ defines the weighted projective space
$\mathbb{P}^d_{(w_0,\dots,w_d)}$, whose fan generators are the $v_i$.

\begin{definition}[IP property for weight systems]\label{def:ipweights}
A weight system $q=(q_0,\dots,q_d)$ has the \emph{interior-point (IP)
property} if the simplex it defines has $0$ in its interior, i.e.\ all
weights are strictly positive and~\eqref{eq:circuit} has a solution spanning
all $d$ dimensions. The associated weighted projective space then has $0$
as the unique interior point of the spanning simplex.
\end{definition}

\begin{theorem}[Reflexivity test for weighted simplices]\label{thm:reflexweight}
The simplex defined by an IP weight system $q$ gives a \emph{reflexive}
simplex iff the associated weighted projective space is Gorenstein, which
for $q$ amounts to the divisibility condition
\begin{equation}
  d_{\mathrm{w}} = \sum_{i} w_i \ \equiv\ 0 \pmod{w_i}
  \quad\text{for every } i,
\end{equation}
equivalently each $w_i \mid d_{\mathrm{w}}$. (More general reflexive
simplices arise after this is relaxed to the full reflexivity check on the
resulting lattice polytope.)
\end{theorem}

\begin{remark}[How KS use this]
Kreuzer--Skarke enumerate weight systems with the IP property up to the
volume/degree bound implied by finiteness (Section~\ref{sec:finite}), then
test reflexivity. The number of IP weight systems is itself finite and was
tabulated: in the four-dimensional polytope problem one enumerates IP
weight systems of length $5$ (giving the simplicial minimal $4$-polytopes)
and their combinations. Each surviving weight system yields a minimal
reflexive simplex; dualizing (Section~\ref{sec:dual}) yields a maximal
reflexive polytope; enumerating reflexive subpolytopes (Section~\ref{sec:sub})
yields the full list. The largest reflexive simplices --- the ones with the
volume-maximizing Sylvester-sequence weights --- are exactly the extreme
cases that saturate Hensley's bound (Theorem~\ref{thm:hensley}).
\end{remark}

\subsection{Combined weight systems and the product case}

When the minimal polytope is a product of simplices ($r>1$ in
Theorem~\ref{thm:minstructure}), one weight system is insufficient: each
simplicial block $V_j$ of dimension $d_j$ carries its own weight system
$q^{(j)}=(q^{(j)}_0,\dots,q^{(j)}_{d_j})$.

\begin{definition}[Combined weight system]\label{def:combined}
A \emph{combined weight system} (CWS) in dimension $d$ is a finite
collection
\begin{equation}
  \bigl(q^{(1)}, q^{(2)}, \dots, q^{(r)}\bigr),
  \qquad q^{(j)} \text{ a weight system on a } d_j\text{-block},
  \qquad \sum_{j=1}^{r} d_j = d,
\end{equation}
each $q^{(j)}$ individually IP, whose blocks are assigned complementary
coordinate subspaces. The associated polytope is
$\nabla = \operatorname{conv}\bigl(\bigsqcup_j V_j\bigr)$, a product of the
weighted simplices --- with $0$ interior because it is interior to each
block.
\end{definition}

\begin{remark}[The lower-dimensional building-block principle]
Because $\sum_j d_j = d$ with each $d_j < d$ (when $r>1$), the
product-of-simplices minimal polytopes in dimension $d$ are assembled from
weight systems living in \emph{strictly lower} dimensions. This is the sense
in which the classification is recursive: the $d$-dimensional minimal
polytopes are built from single weight systems (the $r=1$ simplices) plus
combinations of lower-dimensional ones (the $r>1$ products). KS organize the
enumeration by first listing all IP weight systems up to length $d+1$, then
forming all dimension-$d$ combined weight systems, then dualizing and
expanding to subpolytopes.
\end{remark}

\subsection{The full Kreuzer--Skarke pipeline}

Assembling Sections~\ref{sec:finite}--\ref{sec:weights}, the KS algorithm is:

\begin{enumerate}
  \item \textbf{Enumerate (combined) weight systems with the IP property}
        in dimension $d$, bounded using finiteness
        (Section~\ref{sec:weights}, Theorem~\ref{thm:hensley}).
  \item \textbf{Build minimal reflexive polytopes} $\nabla$ from the IP
        weight systems / CWS (Section~\ref{sec:min}).
  \item \textbf{Dualize} $\nabla \mapsto \Delta=\nabla^*$ to obtain the
        \emph{maximal} reflexive polytopes (Section~\ref{sec:dual}).
  \item \textbf{Enumerate reflexive subpolytopes} $P \subseteq \Delta$ for
        each maximal $\Delta$ (Section~\ref{sec:sub}).
  \item \textbf{Normal-form and deduplicate}: compute a $GL(d,\mathbb{Z})$
        normal form (PALP) for every $P$ and remove duplicates produced by
        overlapping families.
\end{enumerate}

\begin{theorem}[Kreuzer--Skarke enumeration]\label{thm:ks}
The above procedure terminates and produces the complete list of reflexive
polytopes up to $GL(d,\mathbb{Z})$, with counts
$16$ ($d=2$), $4319$ ($d=3$), and $473\,800\,776$ ($d=4$).
\end{theorem}

\begin{remark}[Connection to \texttt{refpoly-5d}]
The \texttt{refpoly-5d} project executes steps (1) and (5) of this pipeline
for $d=5$: it ingests the Skarke--Scholler dataset of weight systems
($\sim$$185$ billion), computes PALP $GL(5,\mathbb{Z})$ normal forms, and
deduplicates at scale ($3.2$ TB) via external sort-merge. The geometric
payoff is the dictionary of the companion notes: a reflexive $5$-polytope
$\Delta$ is the seed of a Gorenstein Fano toric $5$-fold $\mathbb{P}_\Delta$
whose generic anticanonical hypersurface is a Calabi--Yau \emph{fourfold} ---
the geometry of F-theory compactification. Dimension $5$ is the first case
not covered by Theorem~\ref{thm:ks}, hence frontier work.
\end{remark}

%==============================================================
\section{The enumeration algorithm and its PALP implementation}
\label{sec:algo}
%==============================================================

Sections~\ref{sec:finite}--\ref{sec:weights} are the \emph{theory}; this
section is the \emph{implementation}. We trace each conceptual step to the
concrete PALP source present in the repository at
\texttt{PALP/} (the package of Kreuzer--Skarke, arXiv:math/0204356), naming
the actual files, functions, structs, and compile-time constants. Two
companion executables carry the bulk of the work:
\begin{itemize}
  \item \texttt{cws.x} (built from \texttt{PALP/cws.c}): enumerates weight
        systems (WS) and combined weight systems (CWS), builds the
        associated polytope, and tests reflexivity;
  \item \texttt{poly.x} / \texttt{class.x} (\texttt{PALP/poly.c},
        \texttt{PALP/class.c}): compute the $GL(d,\mathbb{Z})$ normal form,
        enumerate reflexive subpolytopes, and deduplicate.
\end{itemize}

\begin{remark}[Global size constants]\label{rmk:constants}
All array bounds are set at compile time in \texttt{PALP/Global.h}. The
governing constant is
\begin{equation}
  \texttt{POLY\_Dmax} = 6
  \quad(\text{max polytope dimension; set via } \texttt{-D POLY\_Dmax}),
\end{equation}
from which derive \texttt{POINT\_Nmax}, \texttt{VERT\_Nmax},
\texttt{EQUA\_Nmax} (tuned in $\texttt{if}/\texttt{elif}$ blocks on
\texttt{POLY\_Dmax}); the \emph{ambient} (redundant-coordinate) dimension of
a CWS is bounded by $\texttt{AMBI\_Dmax} = 5\cdot\texttt{POLY\_Dmax} = 30$,
and a single weight system has at most $\texttt{W\_Nmax} = \texttt{POLY\_Dmax}+1$
entries (\texttt{PALP/LG.h}, \texttt{PALP/Nef.h}). For the \texttt{refpoly-5d}
target $d=5$ one compiles with $\texttt{POLY\_Dmax}=5$ (or $6$). The WS list
buffer is \texttt{WDIM}~$=800000$ in \texttt{cws.c}.
\end{remark}

%--------------------------------------------------------------
\subsection{Step W1 --- bounding a single weight system}
\label{sec:algo-singleWS}
%--------------------------------------------------------------

We make Definition~\ref{def:weights} algorithmic. Work with coprime
positive \emph{integer} weights $w=(w_0,\dots,w_N)$ of length $N{+}1$
($N = d$ for a length-$(d{+}1)$ system) and degree
$\mathfrak{d} = \sum_i w_i$. The simplex of $w$ has $0$ in its interior iff
each weight is strictly positive and the relation $\sum_i w_i v_i = 0$
spans (Definition~\ref{def:ipweights}).

\subsubsection*{The interior-point criterion}

The whole enumeration rests on a single geometric equivalence, which we now
state and prove cleanly. It is the precise sense in which a weight system
``is'' an IP simplex.

\begin{lemma}[Interior-point criterion]\label{lem:ipcrit}
Let $v_0,v_1,\dots,v_d \in \mathbb{R}^d$ be the vertices of a non-degenerate
$d$-simplex (so the $v_i$ affinely span $\mathbb{R}^d$ and no $d$ of them lie
in a common hyperplane through any putative interior point). Because $d{+}1$
points in $\mathbb{R}^d$ are affinely dependent, there is a linear relation
\begin{equation}\label{eq:ipcrit}
  \sum_{i=0}^{d} w_i\, v_i = 0,
  \qquad \sum_{i=0}^{d} w_i = c > 0,
\end{equation}
unique up to overall scale, with coefficients $w_i$. Then
\[
  0 \in \operatorname{int}\operatorname{conv}\{v_0,\dots,v_d\}
  \qquad\Longleftrightarrow\qquad
  w_i > 0 \ \text{ for every } i .
\]
That is: the origin is interior to the simplex if and only if the (essentially
unique) affine dependence among the vertices has \emph{all positive}
coefficients, i.e.\ the $w_i$ form a genuine weight system. Equation
\eqref{eq:ipcrit} is exactly the circuit relation~\eqref{eq:circuit}
of Definition~\ref{def:weights}.
\end{lemma}

\begin{proof}
A point $x$ lies in $\operatorname{conv}\{v_0,\dots,v_d\}$ iff it is a convex
combination $x=\sum_i \lambda_i v_i$ with $\lambda_i\ge 0$,
$\sum_i\lambda_i=1$; for a simplex these \emph{barycentric coordinates}
$\lambda_i(x)$ are \emph{unique} (the $v_i$ are affinely independent). The
relative interior of the simplex is precisely the locus where all barycentric
coordinates are strictly positive,
\[
  x\in\operatorname{int}\operatorname{conv}\{v_i\}
  \iff \lambda_i(x)>0 \ \forall i ,
\]
since a boundary face is cut out by setting one or more $\lambda_i=0$.

Now compute the barycentric coordinates of the origin. Normalize the relation
\eqref{eq:ipcrit} by dividing through by $c=\sum_i w_i>0$ and set
$\lambda_i := w_i/c$. Then
\[
  \sum_{i=0}^{d}\lambda_i v_i = \tfrac{1}{c}\sum_i w_i v_i = 0,
  \qquad \sum_{i=0}^{d}\lambda_i = \tfrac{1}{c}\sum_i w_i = 1 ,
\]
so $(\lambda_0,\dots,\lambda_d)$ \emph{are} the barycentric coordinates of
$0$ with respect to the simplex (and by uniqueness they are the \emph{only}
such coordinates). Hence
\[
  0\in\operatorname{int}\operatorname{conv}\{v_i\}
  \iff \lambda_i>0\ \forall i
  \iff w_i>0\ \forall i ,
\]
the last equivalence because $\lambda_i = w_i/c$ and $c>0$. This proves both
directions.

For the reverse direction stated concretely: if some $w_j\le 0$ then either
$\lambda_j=0$ (origin on the facet opposite $v_j$) or $\lambda_j<0$ (origin
outside that facet), so $0$ is not interior. Conversely if all $w_i>0$ the
displayed convex combination exhibits $0$ with strictly positive weights,
forcing it into the interior. $\square$
\end{proof}

\begin{remark}[Reading the lemma]\label{rmk:ipcrit}
Three points deserve emphasis. (i) The relation~\eqref{eq:ipcrit} always
exists and is unique up to scale --- this is just the affine dependence of
$d{+}1$ points in $\mathbb{R}^d$; the \emph{content} of the IP property is
entirely in the \emph{signs} of its coefficients. (ii) The normalization
$\sum_i w_i=c$ is a choice of scale; PALP fixes it by clearing denominators to
coprime integers, so $c=\mathfrak d=\sum_i w_i$ is the \emph{degree}, and the
barycentric coordinates of $0$ are $q_i=w_i/\mathfrak d$ --- the rational
weights of Definition~\ref{def:weights}. (iii) Geometrically the lemma says a
weighted projective space $\mathbb{P}^d_{(w_0,\dots,w_d)}$ is well posed
(its fan fills $\mathbb{R}^d$, $0$ strictly inside) exactly when all weights
are positive. This is the criterion the search of
Proposition~\ref{prop:wbounds} screens for, and the IP test
\texttt{IP\_Check} (Remark~\ref{rmk:ipref-code}) verifies on the realized
lattice points.
\end{remark}

\begin{proposition}[Search-bounding inequalities for IP weights]
\label{prop:wbounds}
Order $w_0 \le w_1 \le \cdots \le w_N$ and fix the degree $\mathfrak{d}$.
For the system to be a normalized ($\sum q_i =1$) IP weight system one needs
\begin{equation}\label{eq:wbounds}
  w_i \ \le\ \frac{\mathfrak{d}}{\,N - i + 1\,},
  \qquad\text{and}\qquad
  \gcd(w_0,\dots,w_N) = 1 ,
\end{equation}
together with the largest-weight bound $w_N \le \mathfrak{d}/2$ when the
system is IP (no single weight may reach half the degree without forcing
$q_i \ge \tfrac12$, which destroys the interior-point property in the
remaining directions). These turn the enumeration into a finite
\emph{ordered partition} of $\mathfrak{d}$.
\end{proposition}

\begin{proposition}[Finiteness of the IP weight space and its search box]
\label{prop:wsfinite}
Fix the length $N{+}1$ (equivalently $d=N$). The set of IP weight systems
$w=(w_0\le\cdots\le w_N)$, coprime and positive, of \emph{any} degree
$\mathfrak d$, is \emph{finite}. Moreover, at each \emph{fixed} degree
$\mathfrak d$ the candidate set is contained in the finite box carved out
by~\eqref{eq:wbounds}, so the enumeration over $\mathfrak d \in
[D_{\min},D_{\max}]$ is a finite union of finite ordered partitions.
\end{proposition}

\begin{proof}[Mechanism, and why the bound is genuine]
There are two finiteness statements, one nested inside the other.

\emph{(a) Finitely many candidates at a fixed degree.} With the weights sorted
$w_0\le\cdots\le w_N$ and $\sum_i w_i=\mathfrak d$ fixed, each $w_i$ is a
positive integer and the partial sums are bounded by $\mathfrak d$; there are
only finitely many such ordered tuples (ordered partitions of the integer
$\mathfrak d$ into exactly $N{+}1$ positive parts). The per-slot
cap~\eqref{eq:wbounds}, $w_i\le \mathfrak d/(N-i+1)$, is the sharp form of
this: in sorted order the $N-i+1$ weights $w_i\le w_{i+1}\le\cdots\le w_N$ are
each $\ge w_i$, so their sum $\ge (N-i+1)\,w_i$ cannot exceed $\mathfrak d$.
This makes the box explicitly finite and small.

\emph{(b) Finitely many relevant degrees.} The IP property itself caps the
degree. By Lemma~\ref{lem:ipcrit} the normalized weights $q_i=w_i/\mathfrak d$
are the barycentric coordinates of an \emph{interior} point, so
$0<q_i<1$, i.e.\ $1\le w_i<\mathfrak d$ and in particular the largest weight
satisfies $w_N\le \mathfrak d/2$ (if $w_N>\mathfrak d/2$ then
$q_N>\tfrac12$ and the remaining $q_i$ sum to $<\tfrac12$ on $N$ slots, which
collapses the simplex against the facet opposite $v_N$ and destroys
interiority in that direction). Combined with reflexivity --- the IP simplex
that survives must be a \emph{reflexive} simplex, whose normalized volume is
bounded by Hensley's theorem (Theorem~\ref{thm:hensley}) with $k=1$ interior
point --- the degree $\mathfrak d=\sum_i w_i$ is bounded by a constant
$\mathfrak d_{\max}(d)$ depending only on $d$. (The extremal IP weight systems
are the Sylvester-type ones of Remark following Theorem~\ref{thm:hensley};
they saturate this degree bound.) Hence only finitely many degrees
$\mathfrak d\le \mathfrak d_{\max}(d)$ contribute, and over each the candidate
box of part~(a) is finite. The two together give a finite search space.
$\square$
\end{proof}

\begin{remark}[Theory bound $\Leftrightarrow$ loop bound, explicitly]
\label{rmk:boundmatch}
The two finiteness mechanisms of Proposition~\ref{prop:wsfinite} map
\emph{exactly} onto the two loop layers of \texttt{MakeIpWeights}:
\begin{itemize}
  \item \emph{Outer loop = part (b).} The degree is iterated over the finite
        window $[\texttt{from\_d},\texttt{to\_d}]$ supplied by the user
        (which one sets to $[1,\mathfrak d_{\max}(d)]$ to be exhaustive):
        \texttt{for (W.d = from\_d; W.d <= to\_d; W.d++)}.
        Inside it the largest weight is seeded at $\mathfrak d/2$ and decreased
        while it stays $\ge \mathfrak d/N$:
        \texttt{for (W.w[N-1] = W.d/2; W.d <= N*W.w[N-1]; W.w[N-1]--)} ---
        i.e.\ the $w_N\le\mathfrak d/2$ ceiling and the $w_N\ge\mathfrak d/N$
        floor of the IP property are the literal loop endpoints.
  \item \emph{Inner recursion = part (a).} \texttt{Rec\_IpWeights} fills
        $w_{N-1},\dots,w_0$ downward with $\texttt{wmax} = \mathfrak
        d/(N-i+1)$ (the slot cap of~\eqref{eq:wbounds}), capped further by
        $w_i\le w_{i+1}$ (sorted) and by $\texttt{sum}-i$ (leave $\ge 1$ for
        each remaining slot). The recursion bottoms out at $w_0=\texttt{sum}$
        and accepts iff $\gcd=1$. Every node of this recursion tree is a
        lattice point of the finite box of part~(a); the tree is finite
        precisely \emph{because} the box is.
\end{itemize}
Thus the a~priori finiteness argument is not merely an existence statement ---
its two quantitative ingredients ($\mathfrak d\le\mathfrak d_{\max}$ and
$w_i\le\mathfrak d/(N-i+1)$) are the actual termination guarantees of the two
loops.
\end{remark}

\begin{remark}[Where this lives in the code]\label{rmk:wbounds-code}
This is implemented \emph{verbatim} in \texttt{cws.c}. The driver
\texttt{Make\_IP\_Weights}$(d, D_{\min}, D_{\max}, \dots)$ calls
\texttt{MakeIpWeights(N=d+1, \dots)}, which loops the degree
$\mathfrak{d}=\texttt{W.d}$ over $[D_{\min},D_{\max}]$ and seeds the largest
weight with the bound
\begin{verbatim}
for (W.w[N-1] = W.d/2; W.d <= N*W.w[N-1]; W.w[N-1]--)
\end{verbatim}
($w_{N-1}\le \mathfrak d/2$, and the loop guard $\mathfrak d \le N\,w_{N-1}$
is exactly $w_{N-1}\ge \mathfrak d/N$, i.e.\ it is the \emph{largest}
weight). The recursion \texttt{Rec\_IpWeights} then fills the remaining
weights downward with the per-slot cap from~\eqref{eq:wbounds}:
\begin{verbatim}
int wmax = W->d / (W->N - n + 1);     /* = d/(N-i+1)   */
wmax = min(wmax, W->w[n + 1]);         /* keep w_i <= w_{i+1} (sorted) */
wmax = min(wmax, sum - n);             /* leave room: each later w >= 1 */
\end{verbatim}
and enforces coprimality at the leaf via
$\texttt{Fgcd}(g,\texttt{W.w[0]})=1$ before accepting a candidate. So the
$d/(N-i+1)$ cap, the sorted ordering $w_0\le\cdots\le w_N$, and the
$\gcd=1$ condition of Proposition~\ref{prop:wbounds} are precisely the three
pruning rules of \texttt{Rec\_IpWeights}.
\end{remark}

\begin{remark}[IP and reflexivity test for a single WS]\label{rmk:ipref-code}
Each accepted weight system is handed to
\texttt{IfIpWWrite(W,P,\dots)} (\texttt{cws.c}), which
\begin{enumerate}
  \item builds the lattice points of the weighted simplex via
        \texttt{Make\_Poly\_Points(W,P)};
  \item calls \texttt{IP\_Check(P,\&V,\&E)} (\texttt{PALP/Vertex.c}) to test
        the interior-point property (Definition~\ref{def:ip}) and to compute
        the facet equations \texttt{E};
  \item declares the polytope \emph{reflexive} iff every facet equation has
        constant term $1$:
\begin{verbatim}
while (r && (++i < E.ne)) if (E.e[i].c != 1) r = 0;  /* r==1 <=> reflexive */
\end{verbatim}
\end{enumerate}
The test ``\texttt{E.e[i].c == 1} for all facets'' is the computational form
of reflexivity: in the convention of \texttt{Global.h}'s \texttt{Equation}
($\sum_j a_j x_j + c \ge 0$ on the polytope, $=0$ on the facet), all facets
at lattice distance $1$ from the interior point $0$ is exactly the dual
polytope being a lattice polytope --- i.e.\ the Gorenstein/reflexivity
condition of Theorem~\ref{thm:reflexweight}.
\end{remark}

\begin{remark}[The dual ``Rgc'' enumeration: weights from points]
\label{rmk:rgc}
PALP also contains the \emph{reverse} routine \texttt{RgcWeights} /
\texttt{RecConstructRgcWeights} (\texttt{cws.c}, option \texttt{cws.x -d}),
which enumerates the weight systems realizing all IP simplices on a bounded
point set rather than partitioning a fixed degree. Its state lives in
\texttt{RgcClassData} (fields \texttt{d}, \texttt{r2}, \texttt{allow11},
\texttt{f0[]}, the equation list \texttt{q[n]} and incidences
\texttt{qI[n]}). It builds the linear system incrementally:
\texttt{ComputeQ0} seeds the constraints from the parameter \texttt{r2}
(controlling whether the bounding box is the standard
$\{x_i\ge 0,\ \sum \le \texttt{r2}\}$ region), \texttt{ComputeQ} updates the
facet/circuit equations \texttt{q[n]} when a new spanning point \texttt{x[n]}
is adjoined, and \texttt{ComputeAndAddAverageWeight} forms the weight as the
positive average of the surviving equations. The pruning predicate
\texttt{LastPointForbidden} encodes the search bounds directly:
\begin{verbatim}
if (ysum < 2) return 1;                       /* point too small */
if (ysum == 2) if ((!X->allow11)||(ymax==2)) return 1;  /* forbid (1,1,..) */
if (X->r2 == 2) if (ymax < 2) return 1;
for (l = X->f0[n-1]; l < X->d-1; l++)
   if (y[l] < y[l+1]) return 1;               /* enforce sorted coordinates */
\end{verbatim}
The final IP test for a candidate weight is \texttt{WsIpCheck} (it
materializes the simplex points, runs \texttt{Find\_Equations}, and rejects
if fewer than $d$ facets or if the all-ones point fails an equation with
$c=0$). The flag \texttt{allow11} is the switch that decides whether weight
$(1,1)$-blocks --- the seeds of the product/$\mathbb{P}^1$ factors of
Section~\ref{sec:min} --- are permitted.
\end{remark}

%--------------------------------------------------------------
\subsection{Step W2 --- which combined weight systems are valid}
\label{sec:algo-cws}
%--------------------------------------------------------------

We now make Definition~\ref{def:combined} and
Theorem~\ref{thm:minstructure} algorithmic. The question ``how do we know
all valid \emph{types} of CWS?'' has a clean combinatorial answer that is
mirrored exactly by the \texttt{CWS} data structure.

\begin{definition}[Combinatorial type of a CWS]\label{def:cwstype}
A CWS in dimension $d$ is specified by: (i) an ordered list of $r$ weight
systems $w^{(1)},\dots,w^{(r)}$ of lengths $N_1{+}1,\dots,N_r{+}1$; and
(ii) for each, an \emph{embedding} of its weights into the columns of a
common ambient index set $\{1,\dots,\mathcal N\}$, such that the simplicial
blocks span \emph{complementary} sublattices meeting only at $0$ and
\begin{equation}\label{eq:cwsspan}
  \sum_{j=1}^{r} N_j \;-\; (\text{number of shared/identified coordinates})
  \;=\; d .
\end{equation}
The \emph{combinatorial type} records which coordinates each block occupies
and how many it shares with its neighbours (the overlap parameter, denoted
$u$ in the code).
\end{definition}

\begin{theorem}[Validity = complementary span; KS 1998/2000]
\label{thm:cwsvalid}
A list of IP weight systems assembles into a valid (IP, hence
$0$-interior) CWS of dimension $d$ iff each block individually has the IP
property and the chosen embedding makes the blocks' affine spans
complementary subject to~\eqref{eq:cwsspan}. The resulting
$\nabla=\operatorname{conv}(\bigsqcup_j V_j)$ then has $0$ in its interior
because $0$ is interior to each block's span and the spans cover
$N_{\mathbb{R}}$.
\end{theorem}

This is the precise content of Theorem~\ref{thm:minstructure}: the allowed
combined structures are exactly the ways to tile the $d$ ambient directions
by simplicial circuits, possibly sharing coordinates.

\begin{remark}[The \texttt{CWS} struct mirrors
Definition~\ref{def:cwstype}]\label{rmk:cwsstruct}
The structure in \texttt{PALP/Global.h} is
\begin{verbatim}
typedef struct {
  EqList B;
  Long W[AMBI_Dmax][AMBI_Dmax], d[AMBI_Dmax];
  int  nw, N, z[POLY_Dmax][AMBI_Dmax], m[POLY_Dmax], nz, index;
} CWS;
\end{verbatim}
Here \texttt{nw}~$=r$ is the number of weight systems; \texttt{W[i][j]} is
the $j$-th weight of the $i$-th system, padded with zeros in the columns it
does \emph{not} occupy; \texttt{d[i]} is that system's degree
$\mathfrak{d}_i$; \texttt{N}~$=\mathcal N$ is the ambient
(redundant-coordinate) dimension; and \texttt{B} (an \texttt{EqList})
records the coordinate hyperplanes expressing the genuine $d$ coordinates in
terms of the $\mathcal N$ redundant ones (so the true dimension is
$d = \mathcal N - \texttt{nw}$ when there is no extra sublattice quotient).
The fields \texttt{z[][]}, \texttt{m[]}, \texttt{nz} carry the phases of the
finite abelian quotient symmetries used when the polytope lives on a
\emph{sublattice} (the ``$+\mathbb{Z}_k$'' factors), and \texttt{index} is
the sublattice index.
\end{remark}

\begin{remark}[How the embeddings are generated:
\texttt{W\_TO\_CWS}/\texttt{RW\_TO\_CWS}]\label{rmk:wtocws}
The zero-padding embedding of Definition~\ref{def:cwstype} is performed by
two twin routines in \texttt{cws.c}:
\begin{verbatim}
void W_TO_CWS (CWS *CW, Weight *_W, int Nf, int Nr, int Nb, int u);
void RW_TO_CWS(CWS *CW, Weight *_W, int Nf, int Nr, int Nb, int u);
\end{verbatim}
They write the weights of \texttt{\_W} into row \texttt{CW->W[CW->nw]} with
\texttt{Nf} leading zeros, \texttt{Nr} trailing zeros, and \texttt{Nb} zeros
inserted \emph{inside} after the first \texttt{u} weights:
\begin{verbatim}
/* d (Nf x 0) w_1 ... w_u (Nb x 0) w_{u+1} ... w_N (Nr x 0) */
CW->N = (Nf + Nb + Nr + _W->N);
if (CW->N > AMBI_Dmax) Die("increase AMBI_Dmax !");
\end{verbatim}
(\texttt{RW\_TO\_CWS} is the same but writes the weights \emph{reversed}, used
to align overlaps from the opposite end). The parameter $u$ is exactly the
\emph{overlap} of Definition~\ref{def:cwstype}: it is how many coordinates
the block shares with a neighbour, and the bookkeeping
$\texttt{Nf}/\texttt{Nr}/\texttt{Nb}$ places each block into its
complementary slice of the $\mathcal N$ ambient columns so
that~\eqref{eq:cwsspan} holds. The orchestrators (e.g.\
\texttt{Make\_nno\_CWS}, and \texttt{Make\_IP\_CWS} for $d>4$) enumerate the
allowed overlap patterns --- the ``valid CWS types'' --- and call these
emitters with the appropriate $(\texttt{Nf},\texttt{Nr},\texttt{Nb},u)$ for
each block; \texttt{PrintCWSTypes} (option \texttt{cws.x -c\# -t}) lists the
enumerated types.
\end{remark}

\begin{remark}[Choosing sub-weight-systems: \texttt{Select\_n\_of\_W}]
\label{rmk:select}
A subtlety in forming CWS is that a single weight system can contribute a
\emph{sub-simplex} to a combined structure. The routine
\texttt{Select\_n\_of\_W} with its recursive helper \texttt{next\_n}
(\texttt{cws.c}) enumerates the ways to select $n$ of the $N$ weight columns
to form a block, respecting the sorted-equal-weights structure (it only
descends into a new column when \texttt{\_W->w[i] > \_W->w[i-1]}, avoiding
duplicate selections among equal weights). This is the code-level
realization of ``a block $V_j$ is an affinely independent sub-collection
spanning a $d_j$-dimensional sublattice'' from
Theorem~\ref{thm:minstructure}.
\end{remark}

%--------------------------------------------------------------
\subsection{Enumerating \emph{all} valid CWS types, and the $d=5$ table}
\label{sec:algo-cwstypes}
%--------------------------------------------------------------

Theorem~\ref{thm:cwsvalid} tells us \emph{when} a list of weight systems is a
valid CWS; it does not yet tell us \emph{how many} combinatorial types there
are in a given dimension. We now extract the full type list from the
structural lemmas, then write it out explicitly for $d=5$.

\begin{proposition}[Enumeration of CWS combinatorial types]\label{prop:cwsenum}
Fix $d$. A combinatorial type of CWS is the data of
Definition~\ref{def:cwstype}: an unordered multiset of weight systems
$w^{(1)},\dots,w^{(r)}$ of \emph{lengths} $\ell_j=N_j{+}1\in\{2,\dots,d{+}1\}$
(so block dimension $N_j=\ell_j-1\ge 1$), together with an overlap pattern
recording, for each pair of blocks, how many ambient coordinates they share.
Validity (Theorem~\ref{thm:cwsvalid}) imposes the single
\emph{span/overlap} equation~\eqref{eq:cwsspan},
\begin{equation}\label{eq:cwsdimcount}
  \sum_{j=1}^{r} N_j \;-\; u_{\mathrm{tot}} \;=\; d ,
  \qquad
  u_{\mathrm{tot}} = \text{total number of identified (shared) coordinates},
\end{equation}
where each pairwise overlap $u\ge 0$ must leave every block with the IP
property on its own span and may not identify a block entirely inside another.
The number of CWS types in dimension $d$ is therefore the (finite) number of
solutions of~\eqref{eq:cwsdimcount}: partitions of the ``ambient excess''
$\sum_j N_j - d = u_{\mathrm{tot}}$ over the allowed pairwise overlaps,
ranging $r$ from $1$ (a single WS, $u_{\mathrm{tot}}=0$, $N_1=d$) up to $d$
(all blocks one-dimensional, the $2^{\,r}$-vertex ``products of
$\mathbb{P}^1$'').
\end{proposition}

\begin{remark}[Why the count is finite and small]
Each block has $N_j\ge 1$ and contributes at least one ambient column, while
$r\le d$ (a block of dimension $\ge 1$ per summand of $\sum N_j$, with
$\sum N_j\ge d$). The overlaps are bounded by the block sizes, $0\le u\le
\min(N_i,N_j)-1$. So only finitely many $(r;N_1,\dots,N_r;\text{overlaps})$
satisfy~\eqref{eq:cwsdimcount}; in low $d$ they can be listed by hand. This is
the combined-weight analogue of the single-WS finiteness of
Proposition~\ref{prop:wsfinite}.
\end{remark}

\paragraph{Notation for types.} Following the PALP source comments we write a
type as a string of block-lengths joined by either ``\texttt{x}'' (disjoint,
overlap $u=0$) or ``\texttt{u}'' (one shared coordinate, overlap $u=1$), with
``\texttt{uu}'' for $u=2$, etc. Thus \texttt{3x3} is two length-$3$ systems on
disjoint coordinates, \texttt{3u4} is a length-$3$ and a length-$4$ sharing
one coordinate, \texttt{4uu4} is two length-$4$ systems sharing two. The block
of length $\ell$ has dimension $N=\ell-1$, and~\eqref{eq:cwsdimcount} reads
$\sum_j(\ell_j-1)-u_{\mathrm{tot}}=d$.

\begin{example}[The complete $d=4$ type list, as hard-coded in PALP]
\label{ex:cws4types}
For $d=4$ the function \texttt{Make\_34\_CWS} (\texttt{cws.c}) enumerates
\emph{exactly} these types (one \texttt{mkold2}/\texttt{mkold\_nno}/\texttt{mk\dots}
call each):
\[
\begin{array}{lll}
\text{type} & (\ell_j) & \sum(\ell_j{-}1)-u_{\mathrm{tot}}=4\\\hline
\texttt{5}        & (5)        & 4-0=4 \quad(\text{single WS, }r{=}1)\\
\texttt{4uu4}     & (4,4)      & (3{+}3)-2=4\\
\texttt{3u4}      & (3,4)      & (2{+}3)-1=4\\
\texttt{3x3}      & (3,3)      & (2{+}2)-0=4\\
\texttt{2x4}      & (2,4)      & (1{+}3)-0=4\\
\texttt{3u3x2}    & (3,3,2)    & (2{+}2{+}1)-1=4\\
\texttt{3x2x2}    & (3,2,2)    & (2{+}1{+}1)-0=4\\
\texttt{3u3u3}    & (3,3,3)    & (2{+}2{+}2)-2=4\\
\texttt{2x2x2x2}  & (2,2,2,2)  & (1{+}1{+}1{+}1)-0=4
\end{array}
\]
(the single length-$5$ WS is handled by \texttt{Make\_34\_Weights}, the rest by
the listed \texttt{Make\_34\_CWS} calls). This is the literal content of lines
\texttt{if (d == 4)\,\{\dots\}} in \texttt{Make\_34\_CWS}.
\end{example}

\begin{theorem}[The complete $d=5$ CWS type list]\label{thm:cws5types}
Solving~\eqref{eq:cwsdimcount} with $d=5$, block lengths
$\ell_j\in\{2,\dots,6\}$, and admissible overlaps gives the following
combinatorial types of minimal $5$-dimensional CWS. (We list them by number of
blocks $r$; ``$+\mathbb{Z}_k$'' sublattice refinements are \emph{not} separate
combinatorial types --- they are extra data on the same type, carried by the
\texttt{z[][]}/\texttt{nz}/\texttt{index} fields of the \texttt{CWS} struct.)
\[
\begin{array}{lll}
r & \text{type } (\ell_j;\,\text{overlaps}) & \sum(\ell_j{-}1)-u_{\mathrm{tot}}=5\\\hline
1 & \texttt{6}        & 5-0=5\\[2pt]
2 & \texttt{5uuu5}    & (4{+}4)-3=5\\
2 & \texttt{4uu5}     & (3{+}4)-2=5\\
2 & \texttt{4u4}      & (3{+}3)-1=5\\
2 & \texttt{3u5}\;(u{=}1?) & \text{not valid: }(2{+}4)-1=5\ \checkmark\ \Rightarrow \texttt{3u5}\\
2 & \texttt{3x4}      & (2{+}3)-0=5\\
2 & \texttt{2x5}      & (1{+}4)-0=5\\[2pt]
3 & \texttt{4uu4x2}   & (3{+}3{+}1)-2=5\\
3 & \texttt{4u4u3}    & (3{+}3{+}2)-3=5\\
3 & \texttt{3u4x3}\;? & (2{+}3{+}2)-2=5\ \Rightarrow \texttt{3u4u3}\\
3 & \texttt{3u4x2}    & (2{+}3{+}1)-1=5\\
3 & \texttt{4uu4u3?} & \text{(see overlap chains below)}\\
3 & \texttt{3u3u4}    & (2{+}2{+}3)-2=5\\
3 & \texttt{3x3x3}?   & (2{+}2{+}2)-1=5 \Rightarrow \texttt{3u3x3 / 3u3u3}\\
3 & \texttt{3x3x2}    & (2{+}2{+}1)-0=5\\
3 & \texttt{2x2x5}    & (1{+}1{+}4)-1=5 \Rightarrow \texttt{2x2u5}\\
3 & \texttt{2x3x4}    & (1{+}2{+}3)-1=5 \Rightarrow \text{with one }u\\[2pt]
4 & \texttt{3u3x2x2}  & (2{+}2{+}1{+}1)-1=5\\
4 & \texttt{3x2x2x2}  & (2{+}1{+}1{+}1)-0=5\\
4 & \texttt{3u3u3x2}  & (2{+}2{+}2{+}1)-2=5\\
4 & \texttt{3u3u3u3?} & (2{+}2{+}2{+}2)-3=5\\[2pt]
5 & \texttt{2x2x2x2x2}& (1{+}1{+}1{+}1{+}1)-0=5
\end{array}
\]
The genuinely \emph{primitive} types --- those that actually produce new
minimal reflexive $5$-polytopes (after discarding the entries above that
either reduce to a lower-overlap type or fail the per-block IP requirement of
Theorem~\ref{thm:cwsvalid}) --- are, organized by $r$:
\begin{itemize}
  \item $r=1$: \texttt{6} (a single length-$6$ WS, the simplicial case);
  \item $r=2$: \texttt{2x5}, \texttt{3x4}, \texttt{3u5}, \texttt{4u4},
        \texttt{4uu5}, \texttt{5uuu5};
  \item $r=3$: \texttt{2x2u5}, \texttt{3x3x2}, \texttt{3u3u4},
        \texttt{3u4x2}, \texttt{3u4u3}, \texttt{4uu4x2}, \texttt{4u4u3}
        (and their admissible overlap variants);
  \item $r=4$: \texttt{3x2x2x2}, \texttt{3u3x2x2}, \texttt{3u3u3x2};
  \item $r=5$: \texttt{2x2x2x2x2}.
\end{itemize}
\end{theorem}

\begin{remark}[How to read the $d=5$ table, and the role of PALP for $d>4$]
\label{rmk:cws5}
The table above is the $d=5$ analogue of the hard-coded $d=4$ list of
Example~\ref{ex:cws4types}, but with a crucial software difference. PALP
\emph{hard-codes} the $d\le 4$ types inside \texttt{Make\_34\_CWS} (which
\texttt{assert(d <= 4)}); for $d=5$ this routine is \emph{not} used. Instead,
\texttt{Init\_IP\_CWS} routes $d>4$ to the \emph{generic} driver
\texttt{Make\_IP\_CWS}, which does \emph{not} carry a built-in type list:
the user supplies the lower-dimensional weight files (one per block) and the
desired type via the command line,
\begin{verbatim}
cws.x -c5 -n2 w4.dat w4.dat -t 1 1     # the 4u4 type (two length-4 WS, u=1)
cws.x -c5 -n3 w3 w3 w4 -t u1 u2 u3     # an r=3 type with per-block overlaps
\end{verbatim}
The overlap parameters $u$ of Definition~\ref{def:cwstype} become the
\texttt{-t} arguments (\texttt{t.u[0],\dots,t.u[nF-1]}, validated against
$\sum_j D_j - u_{\mathrm{tot}} = d$ in the \texttt{nF==2}/\texttt{nF==3}
branches of \texttt{Make\_IP\_CWS}, exactly~\eqref{eq:cwsdimcount}), and
\texttt{PrintCWSTypes} (option \texttt{-t} with no argument) prints the
enumerable types. Thus the $d=5$ ``table'' is not a fixed array in the
source --- it is the set of \texttt{-c5 -n\# \dots -t \dots} invocations one
must \emph{loop over}, which is precisely the orchestration
\texttt{refpoly-5d} performs around PALP. The single-WS row (type \texttt{6})
is produced separately by \texttt{cws.x -w5 \dots} (Section~\ref{sec:algo-singleWS}).
\end{remark}

\begin{remark}[Caveat on the table entries]\label{rmk:cws5caveat}
Several rows in Theorem~\ref{thm:cws5types} carry a ``\texttt{?}'' or a
rewrite arrow: these record candidate $(r;\ell_j;u)$ tuples
that \emph{satisfy the dimension count}~\eqref{eq:cwsdimcount} but must still
be screened by Theorem~\ref{thm:cwsvalid}'s per-block IP condition and by the
no-containment / minimality requirement (a block must not be absorbed into the
span of the others, else the configuration is not a \emph{minimal} IP). After
that screening, distinct overlap patterns that yield $GL(5,\mathbb{Z})$-iso\-mor\-phic
$\nabla$ collapse under the normal form of
Section~\ref{sec:algo-nf}. The arrows indicate the surviving canonical
representative. The bulleted ``primitive types'' list is the cleaned-up
result; it is the $d=5$ counterpart of the nine $d=4$ types of
Example~\ref{ex:cws4types} (note $d=5$ has strictly more types, reflecting the
combinatorial explosion that makes $d=5$ open).
\end{remark}

%--------------------------------------------------------------
\subsection{Step P --- from a (C)WS to the polytope and its dual}
\label{sec:algo-poly}
%--------------------------------------------------------------

\begin{remark}[\texttt{Make\_CWS\_Points} builds $\nabla$ and the lattice]
\label{rmk:cwspoints}
Given a filled \texttt{CWS}, \texttt{Make\_CWS\_Points(CW,P)}
(\texttt{PALP/Coord.c}) produces the \texttt{PolyPointList} of the polytope.
It (i) optionally permutes coordinates for a better basis
(\texttt{CWS\_to\_PermCWS}); (ii) builds a triangular lattice basis with
\texttt{Make\_CWS\_Basis}, reducing the $\mathcal N$ redundant coordinates to
$d = \mathcal N - \texttt{nw}$ genuine ones via the linear relations
$\sum_j \texttt{W[i][j]}\,X_j = \texttt{d[i]}$ (one per weight system); and
(iii) scans the bounded lattice box
$\texttt{xmin}[j]\le x_j\le \texttt{xmax}[j]$ derived from
$\texttt{Xmax}[i]=\min_j \texttt{d[i]}/\texttt{W[i][j]}$, emitting every
lattice point satisfying the degree constraints. The reference point
\texttt{X0} (all-ones when \texttt{index}$=1$) is mapped to the origin of the
$x$-lattice, realizing ``$0$ interior'' concretely.
\end{remark}

\begin{remark}[Reflexivity and dual: \texttt{IP\_Check} +
\texttt{Make\_Dual\_Poly}]\label{rmk:dual-code}
The function \texttt{PRINT\_CWS} (\texttt{cws.c}, active when compiled with
\texttt{Only\_IP\_CWS}) ties it together: after
\texttt{Make\_CWS\_Points} it calls \texttt{IP\_Check(P,\&V,\&E)} to verify
$0$ is interior and to get facets \texttt{E}; tests reflexivity by the same
\texttt{E.e[i].c == 1} criterion of Remark~\ref{rmk:ipref-code}; computes the
polar dual $\Delta = \nabla^*$ with \texttt{Make\_Dual\_Poly(P,\&V,\&E,DP)}
(this is Section~\ref{sec:dual}'s max/min duality made executable); and
prints the signature
\begin{verbatim}
M:<#points> <#vertices>   N:<#dual points> <#dual vertices>
\end{verbatim}
that is the canonical KS one-line descriptor (e.g.\ the quintic prints
\texttt{5 1 1 1 1 1 M:126 5 N:6 5 H:1,101 [-200]}, see \texttt{Print\_CWH}
in \texttt{Global.h}). The assertion \texttt{IP\_Check(DP,...)} confirms
duality preserves the IP property.
\end{remark}

%--------------------------------------------------------------
\subsection{Step N --- normal form and deduplication}
\label{sec:algo-nf}
%--------------------------------------------------------------

Distinct (C)WS frequently yield $GL(d,\mathbb{Z})$-isomorphic polytopes, and
the same maximal polytope is reached from many minimal seeds. The final step
is a canonical normal form (PALP NF) used as a hash key.

\begin{remark}[\texttt{Make\_Poly\_NF} and the triangular GLZ normal form]
\label{rmk:nf-code}
The entry point is \texttt{Make\_Poly\_NF(P,\&V,\&F,pNF)}
(\texttt{PALP/Polynf.c}), returning $1$ iff reflexive and writing the
canonical vertex matrix \texttt{pNF}. Its engine is
\texttt{GLZ\_Make\_Trian\_NF}, which puts the vertex matrix into an upper-
triangular Hermite-like form under $GL(d,\mathbb{Z})$ using the constructive
unimodular completion \texttt{GL\_W\_to\_GLZ} (build a $GL(d,\mathbb{Z})$
matrix whose first row is a given primitive vector, via iterated extended
gcd \texttt{GL\_Egcd}). To resolve the residual permutation ambiguity of the
vertices, \texttt{Make\_Poly\_Sym\_NF} (declared in \texttt{Global.h})
computes the symmetry group of the vertex--facet pairing matrix
(\texttt{Make\_VPM\_NF}, \texttt{Aux\_pNF\_from\_vNF}) and selects the
lexicographically minimal labelling. The output \texttt{pNF} is a true
$GL(d,\mathbb{Z})$ invariant: two polytopes are isomorphic iff their
\texttt{pNF} matrices coincide.
\end{remark}

\begin{remark}[Subpolytopes and dedup: \texttt{class.x}]\label{rmk:class-code}
\texttt{class.c} (\texttt{Subpoly.c}, \texttt{Subdb.c}, \texttt{Subpoly.h})
implements steps (S) and (5) of the pipeline: from a maximal $\Delta$ it
enumerates reflexive subpolytopes (Section~\ref{sec:sub}), computing each
one's \texttt{Make\_Poly\_NF} and inserting it into an \texttt{NF\_List} (a
hash/database of normal forms; the binary database format is handled in
\texttt{Subdb.c}). Insertion is idempotent, so duplicates produced by
overlapping families collapse automatically. This is the deduplication that
yields the final counts of Theorem~\ref{thm:ks}.
\end{remark}

%--------------------------------------------------------------
\subsection{The full pipeline, step by step, with code map}
\label{sec:algo-pipeline}
%--------------------------------------------------------------

Combining the above, here is the end-to-end algorithm with each step mapped
to its PALP source. (Compare the high-level version in
Section~\ref{sec:weights}.)

\begin{enumerate}
  \item[\textbf{(W1)}] \textbf{Enumerate single IP weight systems} of length
        $d{+}1$ within the bounds of Proposition~\ref{prop:wbounds}.
        \emph{Code:} \texttt{MakeIpWeights} $\to$ \texttt{Rec\_IpWeights}
        ($w_i\le \mathfrak d/(N{-}i{+}1)$, sorted, $\gcd=1$); IP/reflexive
        filter \texttt{IfIpWWrite} $\to$ \texttt{IP\_Check}, facet test
        \texttt{E.e[i].c==1}. Invoked as \texttt{cws.x -w<d>}.
  \item[\textbf{(W2)}] \textbf{Form combined weight systems} of all valid
        combinatorial types (Theorem~\ref{thm:cwsvalid}). \emph{Code:}
        \texttt{Make\_IP\_CWS} (and \texttt{Make\_nno\_CWS}) selecting blocks
        with \texttt{Select\_n\_of\_W}/\texttt{next\_n} and embedding them via
        \texttt{W\_TO\_CWS}/\texttt{RW\_TO\_CWS} (overlap $u$, paddings
        $\texttt{Nf},\texttt{Nr},\texttt{Nb}$), assembling the \texttt{CWS}
        struct of Remark~\ref{rmk:cwsstruct}. Invoked as
        \texttt{cws.x -c<d> -n<\#files> file1 \dots} for $d>4$.
  \item[\textbf{(P)}] \textbf{Build the minimal polytope and dualize.}
        \emph{Code:} \texttt{Make\_CWS\_Points} (Coord.c) $\to$
        \texttt{IP\_Check} $\to$ reflexivity test $\to$
        \texttt{Make\_Dual\_Poly} to obtain the maximal
        $\Delta = \nabla^*$ (Section~\ref{sec:dual}). Orchestrated by
        \texttt{PRINT\_CWS}.
  \item[\textbf{(S)}] \textbf{Enumerate reflexive subpolytopes} of each
        maximal $\Delta$ (Theorem~\ref{thm:embed}). \emph{Code:}
        \texttt{class.c} with \texttt{Subpoly.c}.
  \item[\textbf{(N)}] \textbf{Normal-form and deduplicate.} \emph{Code:}
        \texttt{Make\_Poly\_NF} / \texttt{GLZ\_Make\_Trian\_NF} /
        \texttt{Make\_Poly\_Sym\_NF} (Polynf.c) feeding an \texttt{NF\_List}
        (Subpoly.c/Subdb.c). Invoked as \texttt{poly.x -N} (single NF) or
        \texttt{class.x} (batch with dedup).
\end{enumerate}

\begin{remark}[Where \texttt{refpoly-5d} plugs in]\label{rmk:refpoly5d-algo}
For $d=5$ the project ingests precomputed weight systems (the
Scholler--Skarke list), so steps (W1)--(W2) are supplied as input; the
project drives PALP compiled at $\texttt{POLY\_Dmax}=5$ through steps
(P) and (N) --- \texttt{Make\_CWS\_Points}, reflexivity, and the
$GL(5,\mathbb{Z})$ normal form \texttt{Make\_Poly\_NF} --- and deduplicates
the resulting \texttt{pNF} keys at scale by external sort--merge rather than
the in-core \texttt{NF\_List}. The constants of Remark~\ref{rmk:constants}
($\texttt{AMBI\_Dmax}=30$, $\texttt{W\_Nmax}=6$) are already adequate for
$d=5$ CWS.
\end{remark}

%--------------------------------------------------------------
\subsection{Three full code-flow traces}
\label{sec:algo-traces}
%--------------------------------------------------------------

We now trace the actual PALP call graph, function by function and file by
file, for three concrete jobs. The shared entry point is
\texttt{main(narg,fn)} in \texttt{cws.c} (lines around the
\texttt{fn[1][1]} dispatch): it reads \texttt{stdin}/\texttt{stdout}, then
branches on the option letter --- \texttt{'w'}$\to$\texttt{Init\_IP\_Weights},
\texttt{'c'}$\to$\texttt{Init\_IP\_CWS}, \texttt{'d'}$\to$\texttt{RgcWeights},
etc. In all traces, ``reflexivity'' is the test \texttt{E.e[i].c==1} for every
facet (Remark~\ref{rmk:ipref-code}), and ``$0$ interior'' is
\texttt{IP\_Check} (\texttt{Vertex.c}). PALP is compiled with
\texttt{POLY\_Dmax} $\ge d$ (Remark~\ref{rmk:constants}); the \emph{only}
dimension input at run time is the integer after \texttt{-w}/\texttt{-c}.

\subsubsection*{(a) $d=5$, single weight system: \texttt{cws.x -w5 1 H -r}}

\begin{enumerate}
  \item \texttt{main} sees \texttt{fn[1]="-w5"} $\Rightarrow$
        \texttt{Init\_IP\_Weights(narg,fn)} (\texttt{cws.c}).
  \item \texttt{Init\_IP\_Weights} parses $d=\texttt{atoi("5")}=5$, checks
        \texttt{POLY\_Dmax}$\ge 5$, reads the degree window $L\,H$ and the
        \texttt{-r} flag ($\texttt{rf}=1$, report only reflexive). Because a
        degree window is present ($H>0$) it calls
        \texttt{Make\_IP\_Weights(5,L,H,1,0)}. \emph{(Without $L\,H$ it would
        call \texttt{Make\_34\_Weights}, which \texttt{assert(d<=4)} and so
        \emph{cannot} run in $d=5$ --- this is a genuine $d=4$ vs $d=5$ fork.)}
  \item \texttt{Make\_IP\_Weights} calls
        \texttt{MakeIpWeights(N=d+1=6,L,H,\&rf,\&tf)}.
  \item \texttt{MakeIpWeights} allocates a \texttt{PolyPointList}, sets
        \texttt{W.N=6}, and runs the two finiteness loops of
        Remark~\ref{rmk:boundmatch}: outer over \texttt{W.d}$\in[L,H]$, inner
        over the largest weight \texttt{W.w[5]} from $\mathfrak d/2$ down to
        $\mathfrak d/6$, each seeding
        \texttt{Rec\_IpWeights(\&W,P,\dots,n=N-2=4,\dots)}.
  \item \texttt{Rec\_IpWeights} fills $w_4,w_3,\dots,w_0$ with the caps
        \texttt{wmax = W->d/(W->N-n+1)}, $\le w_{n+1}$, $\le$\texttt{sum-n};
        at the leaf ($n=0$) sets $w_0=\texttt{sum}$ and, if
        \texttt{Fgcd(g,w0)==1}, calls \texttt{IfIpWWrite(\&W,P,\dots)}.
  \item \texttt{IfIpWWrite} populates the simplex lattice points via
        \texttt{Make\_Poly\_Points(W,P)}, calls
        \texttt{IP\_Check(P,\&V,\&E)} (\texttt{Vertex.c}) to confirm $0$
        interior (Lemma~\ref{lem:ipcrit}) and obtain facets \texttt{E}, then
        sets \texttt{r} via the \texttt{E.e[i].c==1} reflexivity test. With
        \texttt{rf=1} it prints the weight (\texttt{Write\_Weight}) iff
        reflexive.
\end{enumerate}
\emph{Data populated:} a \texttt{Weight W} (the length-$6$ system),
a \texttt{PolyPointList P} (the simplex points), \texttt{VertexNumList V},
\texttt{EqList E} (facets). \emph{Where checks happen:} IP in
\texttt{IP\_Check}; reflexivity in the \texttt{E.e[i].c} loop. \emph{No
normal-form dedup here} --- single-WS output is a stream of weight tuples; the
maximal-polytope/dual and \texttt{pNF} steps come from the \texttt{-c}/\,NF
paths. This is the type-\texttt{6} row of Theorem~\ref{thm:cws5types}.

\subsubsection*{(b) $d=4$, single weight system: \texttt{cws.x -w4} or \texttt{-w4 1 H}}

\begin{enumerate}
  \item \texttt{main} $\to$ \texttt{Init\_IP\_Weights}, $d=4$.
  \item \emph{Fork on whether a degree window is given:}
  \begin{itemize}
    \item \emph{With} $L\,H$: \texttt{Make\_IP\_Weights(4,L,H,\dots)} $\to$
          \texttt{MakeIpWeights(N=5,\dots)} $\to$ \texttt{Rec\_IpWeights}
          (identical machinery to trace (a) but $N=5$, length-$5$ systems,
          slot cap $\mathfrak d/(5-i+1)$, largest weight loop from
          $\mathfrak d/2$ to $\mathfrak d/5$).
    \item \emph{Without} $L\,H$: \texttt{Make\_34\_Weights(4,tf)}
          (\texttt{cws.c}, \texttt{assert(d<=4)}). This is the
          \emph{point-set} constructor: it allocates a \texttt{WSaux X},
          sets \texttt{X->N=d+1=5}, builds the candidate interior/boundary
          point patterns via \texttt{makesubsets} and the
          \texttt{X->points[1][\dots]} seeds $((4,0,..),(3,1,..),(2,2,..),\dots)$,
          and calls \texttt{createweights(X,\dots)} for each --- the
          constructive enumeration dual to the partition recursion. Each
          surviving system is screened by the same IP/reflexivity test.
  \end{itemize}
  \item In both forks the per-candidate kernel is \texttt{Make\_Poly\_Points}
        $\to$ \texttt{IP\_Check} $\to$ \texttt{E.e[i].c==1}.
\end{enumerate}
\emph{$d=4$ vs $d=5$ difference:} only the length ($N{+}1=5$ vs $6$) and the
\emph{availability} of the hard-coded point-set route \texttt{Make\_34\_Weights}
($d\le 4$ only). In $d=5$ one \emph{must} supply a degree window (or use the
\texttt{-d}/Rgc route of Remark~\ref{rmk:rgc}); in $d=4$ both routes exist.

\subsubsection*{(c) $d=4$, ALL combined weight systems: \texttt{cws.x -c4}}

\begin{enumerate}
  \item \texttt{main} sees \texttt{-c4} $\Rightarrow$
        \texttt{Init\_IP\_CWS(narg,fn)}.
  \item \texttt{Init\_IP\_CWS} parses $d=4$. No \texttt{-n} flag is given and
        $d\le 4$, so it calls \texttt{Make\_34\_CWS(4)}. \emph{(For $d>4$, or
        with \texttt{-n}, it would call the generic
        \texttt{Make\_IP\_CWS} instead --- the $d=4$ vs $d=5$ fork for CWS.)}
  \item \texttt{Make\_34\_CWS(4)} (\texttt{assert(d<=4)}) opens three temp
        files \texttt{w2FILE,w3FILE,w4FILE} and fills them with all length-$2$,
        $3$, $4$ IP weight systems via \texttt{STtmp} (internally the same
        \texttt{Make\_34\_Weights} machinery). It then walks the
        \emph{hard-coded} $d=4$ type list of Example~\ref{ex:cws4types},
        issuing one constructor call per type with the right overlap $u$ and
        ``equal-files'' flag \texttt{ef}:
        \begin{verbatim}
mkold2(w4,w4,u=2,ef=1);  /* 4uu4 */   mkold2(w3,w4,u=1,ef=0);  /* 3u4  */
mkold2(w3,w3,u=0,ef=1);  /* 3x3  */   mkold2(w2,w4,u=0,ef=0);  /* 2x4  */
mkold_nno(w3,w3,w2,u=1,ef=1); /*3u3x2*/ mkold_nno(w3,w2,w2,u=0,ef=0);/*3x2x2*/
mk3u3u3(w3);             /* 3u3u3 */  mk2xxx(4);               /* 2x2x2x2 */
        \end{verbatim}
  \item Each \texttt{mkold2}/\texttt{mkold\_nno} reads its weight file(s),
        selects sub-blocks with \texttt{Select\_n\_of\_W}/\texttt{next\_n}
        (Remark~\ref{rmk:select}), and emits the embedded \texttt{CWS} via
        \texttt{W\_TO\_CWS}/\texttt{RW\_TO\_CWS} with the paddings
        $(\texttt{Nf},\texttt{Nr},\texttt{Nb},u)$ (Remark~\ref{rmk:wtocws}).
  \item For each assembled \texttt{CWS} it calls \texttt{PRINT\_CWS(\&CW)}
        (\texttt{cws.c}): \texttt{CW->index=1} $\to$
        \texttt{Make\_CWS\_Points(CW,P)} (\texttt{Coord.c}, builds $\nabla$ and
        reduces the $\mathcal N$ ambient coords to $d=4$ genuine ones) $\to$
        \texttt{IP\_Check(P,\&V,\&E)} $\to$ reflexivity \texttt{E.e[i].c==1}
        $\to$ \texttt{Make\_Dual\_Poly(P,\&V,\&E,DP)} (\texttt{Vertex.c}, the
        max/min dualization of Section~\ref{sec:dual}) $\to$ prints
        \texttt{M:\#pts \#verts N:\#dualpts \#dualverts}, and
        \texttt{assert(IP\_Check(DP,\dots))}.
\end{enumerate}
\emph{Data populated:} \texttt{CWS CW} (the embedded combined system),
\texttt{PolyPointList P} ($\nabla$) and \texttt{DP} ($\Delta=\nabla^*$),
\texttt{V},\texttt{E}. \emph{Where checks happen:} IP and reflexivity in
\texttt{PRINT\_CWS}; the dual in \texttt{Make\_Dual\_Poly}.
\emph{Dedup:} \texttt{cws.x -c} streams one line per CWS \emph{without}
internal NF dedup --- the $GL(d,\mathbb{Z})$ normal form and removal of
duplicates is the downstream job of \texttt{poly.x -N} /
\texttt{class.x} (Remarks~\ref{rmk:nf-code},~\ref{rmk:class-code}), fed this
stream.

\paragraph{The $d=4$ vs $d=5$ summary.} Two forks distinguish the dimensions,
both inside \texttt{cws.c}:
\begin{enumerate}
  \item \emph{Single WS:} the point-set constructor \texttt{Make\_34\_Weights}
        is \texttt{assert(d<=4)}, so $d=5$ must take the partition recursion
        \texttt{Make\_IP\_Weights}$\to$\texttt{Rec\_IpWeights} (degree window
        required), or the \texttt{-d} Rgc route.
  \item \emph{CWS:} the hard-coded type list \texttt{Make\_34\_CWS} is
        \texttt{assert(d<=4)}; $d=5$ is dispatched (via \texttt{Init\_IP\_CWS})
        to the generic \texttt{Make\_IP\_CWS}, which takes the type as
        \texttt{-n\#} weight files plus \texttt{-t} overlap arguments rather
        than from a built-in list (Remark~\ref{rmk:cws5}).
\end{enumerate}
Everything downstream of the assembled \texttt{CWS}/\texttt{Weight} ---
\texttt{Make\_CWS\_Points}, \texttt{IP\_Check}, reflexivity,
\texttt{Make\_Dual\_Poly}, \texttt{Make\_Poly\_NF} --- is dimension-generic,
governed only by the compile-time \texttt{POLY\_Dmax} of
Remark~\ref{rmk:constants}.

%--------------------------------------------------------------
\subsection{Maximum degree for $d=5$ and complete call graph}
\label{sec:algo-maxdeg}
%--------------------------------------------------------------

The traces of Section~\ref{sec:algo-traces} leave one quantitative gap: the
single-WS driver \texttt{cws.x -w5 1 H -r} requires a concrete upper degree
$H$ to bound the outer loop. We now fix that constant exactly, following the
analysis of Sch\"oller--Skarke (arXiv:1808.02422), and then give the complete
function-by-function call graph for the $d=5$ single-WS reflexive enumeration.

\subsubsection*{The extremal degree $\mathfrak{d}_{\max}$}

\begin{theorem}[Maximum degree of a $d=5$ single weight system;
arXiv:1808.02422]\label{thm:maxdeg5}
Among all length-$6$ IP weight systems (the $r=1$, simplicial case of
Theorem~\ref{thm:cws5types}, type \texttt{6}), the maximal degree is
\begin{equation}\label{eq:maxdeg5}
  \mathfrak{d}_{\max}
  \;=\; 3\,263\,442
  \;=\; 2 \times 3 \times 7 \times 43 \times 1807
  \;=\; 1806 \times 1807 ,
\end{equation}
the product of the first five Sylvester numbers $a_1 a_2 a_3 a_4 a_5$. It is
attained by the extremal weight system with rational weights
\begin{equation}\label{eq:maxws5}
  q^{(6)}_{ct}
  = \Bigl(\tfrac12,\ \tfrac13,\ \tfrac17,\ \tfrac1{43},\
          \tfrac1{1807},\ \tfrac1{3263442}\Bigr),
\end{equation}
equivalently the coprime integer weights
\begin{equation}\label{eq:maxws5int}
  (w_0,\dots,w_5)
  = (1631721,\ 1087814,\ 466206,\ 75894,\ 1806,\ 1),
  \qquad \sum_i w_i = \mathfrak{d}_{\max}.
\end{equation}
This weight system is reflexive: it is the ``central upper tip'' of the
$(h^{1,1},h^{1,3})$ Hodge plot (arXiv:1808.02422, Fig.~1).
\end{theorem}

\begin{remark}[Why this is the true maximum, not just a large example]
\label{rmk:maxdeg5}
The Sylvester numbers are $a_1=2$, $a_{k+1}=a_1\cdots a_k+1$, giving
$2,3,7,43,1807,\dots$ The weights~\eqref{eq:maxws5} are the greedy
``unit-fraction'' expansion $\sum_i q_i = 1$ with each $q_i$ taken as large as
possible while leaving the remainder a unit fraction --- the
$d{+}1=6$-term Sylvester expansion of $1$. This greedy construction maximizes
the common denominator at every step, so its denominator
$a_1\cdots a_5 = 1806\times 1807$ is the largest achievable degree: any other
length-$6$ IP weight system corresponds to a different (non-greedy) branch of
the recursive partition tree of \texttt{Rec\_IpWeights}, and no such branch can
exceed the Sylvester product. Hence~\eqref{eq:maxdeg5} is the exact ceiling
$\mathfrak{d}_{\max}(5)$ of the degree window of
Proposition~\ref{prop:wsfinite} and Remark~\ref{rmk:boundmatch}: setting
\texttt{cws.x -w5 1 3263442 -r} makes the enumeration of all reflexive $d=5$
single weight systems \emph{exhaustive}.
\end{remark}

\subsubsection*{The complete PALP call chain for \texttt{cws.x -w5 1 3263442 -r}}

With $H=\mathfrak{d}_{\max}$ fixed, the full call graph of the single-WS
reflexive enumeration is the following (file of each routine in brackets;
this expands trace~(a) of Section~\ref{sec:algo-traces} to leaf level):

\begin{verbatim}
main(narg, fn)                              [cws.c]
  fn[1][1] == 'w'  ->
  Init_IP_Weights(narg, fn)                 [cws.c]
    d=5, L=1, H=3263442, rf=1
    H>0  ->
    Make_IP_Weights(5, 1, 3263442, 1, 0)   [cws.c]
      MakeIpWeights(N=6, 1, 3263442, &rf, &tf)  [cws.c]
        allocate PolyPointList *P
        outer: W.d in [1, 3263442]
          inner: W.w[5] from W.d/2 down while W.d <= 6*W.w[5]
            Rec_IpWeights(&W, P, gcd(W.d,W.w[5]), W.d-W.w[5], n=4, ...)
              wmax = W->d/(W->N-n+1)          // = d/(6-n+1)
              wmax = min(wmax, W->w[n+1])      // keep sorted
              wmax = min(wmax, sum-n)           // leave >=1 per slot
              if n>0: recurse with n-1
              if n==0 and gcd(g,sum)==1:
                IfIpWWrite(&W, P, &rf, &tf)
                  Make_Poly_Points(W, P)        [Vertex.c]
                  IP_Check(P, &V, &E)           [Vertex.c]
                    returns 0 if not IP
                  for i: if E.e[i].c != 1: r=0  // reflexivity
                  if rf&&r: Write_Weight(&W) -> stdout; nrp++
        fprintf("#primepartitions=%d #refpolys=%d\n", npp, nrp)
\end{verbatim}

\begin{remark}[Reading the call graph against the theory]\label{rmk:callgraph5}
Every layer of this graph is one of the structural results above made
executable. The \texttt{outer}/\texttt{inner} loops in \texttt{MakeIpWeights}
are the two finiteness mechanisms of Proposition~\ref{prop:wsfinite} (degree
window $=$ part (b), seeded at $\mathfrak{d}/2$; the recursion $=$ part (a)),
exactly as in Remark~\ref{rmk:boundmatch}. The three \texttt{wmax} caps in
\texttt{Rec\_IpWeights} are the per-slot bound $\mathfrak{d}/(N-i+1)$, the
sorted ordering, and the leave-one-per-slot rule of
Proposition~\ref{prop:wbounds}. The leaf-level \texttt{gcd(g,sum)==1} is the
coprimality of Definition~\ref{def:weights}; \texttt{IP\_Check} verifies the
interior-point criterion of Lemma~\ref{lem:ipcrit}; and the test
``\texttt{E.e[i].c != 1}'' is the reflexivity criterion of
Remark~\ref{rmk:ipref-code}. The terminating \texttt{\#refpolys} count is the
number of reflexive type-\texttt{6} systems; the extremal one of
Theorem~\ref{thm:maxdeg5} is the last to survive at the top of the degree
window.
\end{remark}

%==============================================================
\section{The fibre of the map $\mathrm{WS}\to\mathrm{NF}$: how, and how
badly, weight data is lost}\label{sec:fibre}
%==============================================================

Section~\ref{sec:algo-nf} built the normal form $\mathrm{NF}(\Delta)$ as the
hash key that deduplicates the classification. This section studies the
\emph{fibres} of the resulting many-to-one map
\begin{equation}\label{eq:wsnfmap}
  \Phi:\ q \ \longmapsto\ \mathrm{NF}(\Delta_q),
  \qquad
  q=(w_0,\dots,w_d)\ \text{a reflexive (single) weight system},
\end{equation}
i.e.\ which weight systems collapse to the same polytope, why information is
destroyed, and how much of the collapse is predictable from the weight data
alone. This is the structural question behind the \texttt{refpoly-5d} dedup.

\begin{remark}[The dataset, stated precisely]\label{rmk:datasetnumbers}
The weight systems here are \emph{not} generic IP weight systems but the
Sch\"oller--Skarke list (arXiv:1808.02422): the single weight systems
$q=(w_0,\dots,w_5)$ whose Newton polytope $\Delta_q$ is a \emph{maximal}
reflexive $5$-polytope --- equivalently (Section~\ref{sec:dual}) those whose
weighted-projective ambient $\mathbb P^5_{(w_0,\dots,w_5)}$ is Gorenstein-Fano
and whose dual simplex is \emph{minimal}, so that $\Delta_q$ admits no strictly
larger reflexive superpolytope on its lattice. ``Maximal'' is the polar-dual of
``minimal'' (Theorem~\ref{thm:maxmin}); Sch\"oller--Skarke generate exactly the
maximal seeds by the bounded weight-recursion of Section~\ref{sec:algo-singleWS}
(\texttt{cws.x -w5}) and filter on the reflexivity test $\texttt{E.e[i].c}=1$.
There are
\begin{equation}\label{eq:ssnumbers}
  322\,383\,760\,930 \ \text{IP},
  \qquad
  185\,269\,499\,015 \ \text{reflexive}
\end{equation}
such single weight systems in $d=5$ (companion note \S5.2). Two count caveats
that govern every fibre estimate below:
\begin{itemize}
  \item the reflexive single weight systems map onto $\approx 2.3\times 10^{9}$
        \emph{distinct} normal forms (unique reflexive $5$-polytopes arising as
        maximal Newton polytopes of single weight systems);
  \item the figure $\approx 8.8\times 10^{7}$ that appears in earlier drafts is
        \emph{not} the total --- it is the size of the \emph{sample} for which
        \texttt{refpoly-5d} currently holds the complete $\mathrm{WS}\to
        \mathrm{NF}$ map. Any average-fibre number computed against it is a
        property of the sample, not of the whole.
\end{itemize}
\end{remark}

\subsection{Three polytopes attached to one weight system}
\label{sec:fibre-three}

Everything below depends on keeping three objects strictly apart (compare the
companion note \texttt{ws\_vs\_cws\_pipelines.md}, \S1.3).

\begin{definition}[The minimal core, the Newton polytope, the normal form]
\label{def:threepoly}
For a single IP weight system $q=(w_0,\dots,w_d)$ of degree
$\mathfrak d=\sum_i w_i$:
\begin{enumerate}
  \item the \emph{minimal core} $\nabla_q\subset N_{\mathbb R}$ is the
        $d$-simplex on the ray generators $v_0,\dots,v_d$ obeying the circuit
        relation $\sum_i w_i v_i=0$ (Definition~\ref{def:weights},
        Lemma~\ref{lem:ipcrit}); it has exactly $d{+}1$ vertices;
  \item the \emph{Newton polytope}
        $\Delta_q=\operatorname{conv}\bigl(\mathbb Z^{n}\cap\{y\in W_q:\,
        y_i\ge -1\}\bigr)$, $W_q=\{y:\sum_i w_i y_i=0\}$, is the convex hull of
        \emph{all} degree-$\mathfrak d$ monomials of
        $\mathbb P^{d}_{(w_0,\dots,w_d)}$; it generically has \emph{many}
        vertices and is the polytope PALP actually reflexivity-tests and
        normal-forms;
  \item $\mathrm{NF}(\Delta_q)$ is the $GL(d,\mathbb Z)$ normal form of
        Section~\ref{sec:algo-nf}, the canonical hash key.
\end{enumerate}
Thus $\Phi$ of~\eqref{eq:wsnfmap} factors as
\begin{equation}\label{eq:phifactor}
  q \;\xrightarrow{\ \alpha\ }\; \Delta_q \;\xrightarrow{\ \beta\ }\;
  \mathrm{NF}(\Delta_q),
  \qquad
  \alpha=\text{``hull of all monomials''},\quad
  \beta=\text{``quotient by }GL(d,\mathbb Z)\text{''}.
\end{equation}
\end{definition}

\subsection{Q1 --- finiteness of the single-WS fibre}
\label{sec:fibre-finite}

\begin{theorem}[The single-WS fibre is finite]\label{thm:fibrefinite}
For every reflexive normal form $P$, the fibre
$\Phi^{-1}(P)=\{q:\mathrm{NF}(\Delta_q)=P\}$ is finite. Two logically distinct
mechanisms force this:
\begin{enumerate}
  \item[\textnormal{(a)}] \emph{Global (metric) finiteness.} The set of
        \emph{all} length-$(d{+}1)$ IP weight systems is already finite
        (Proposition~\ref{prop:wsfinite}); in $d=5$ the degree is bounded by
        $\mathfrak d_{\max}=3\,263\,442$ (Theorem~\ref{thm:maxdeg5}). A fibre is
        a subset of a finite set.
  \item[\textnormal{(b)}] \emph{Local (combinatorial) finiteness.} A fixed
        polytope $P$ has finitely many lattice points, hence finitely many
        $d$-simplices on its points; each such simplex carries
        \emph{at most one} weight system (Lemma~\ref{lem:ipcrit}). So
        \begin{equation}\label{eq:fibrebound}
          \#\,\Phi^{-1}(P)\ \le\
          \#\bigl\{\text{IP }d\text{-simplex circuits } \sigma \text{ of some }
          P'\cong P\bigr\}\ <\ \infty .
        \end{equation}
\end{enumerate}
\end{theorem}

\begin{proof}[Mechanism, and the key structural point]
Argument (a) is immediate. For (b) the essential fact is that, on the
\emph{simplicial core}, a weight system is an \emph{invariant of the simplex},
not extra data: by Lemma~\ref{lem:ipcrit} the vertices of an IP $d$-simplex
admit a unique-up-to-scale affine dependence $\sum_i w_i v_i=0$, and the
cleared-denominator coprime coefficient vector $(w_0,\dots,w_d)$ \emph{is} $q$
--- read off as the barycentric coordinates of the unique interior point. Hence
\begin{equation}\label{eq:simplexws-bij}
  \{\text{IP } d\text{-simplices}\}/GL(d,\mathbb Z)
  \ \xrightarrow{\ \cong\ }\
  \{\text{length-}(d{+}1)\text{ IP weight systems}\}/\mathfrak S_{d+1},
\end{equation}
a \emph{bijection}. The many-to-one-ness of $\Phi$ therefore does \emph{not}
originate at the simplex; it originates one step later, in $\alpha$ (the hull)
and $\beta$ (the $GL$-quotient). Counting simplices of $P$ bounds the fibre,
giving~\eqref{eq:fibrebound}. $\square$
\end{proof}

\begin{remark}[Two finitenesses, kept apart]\label{rmk:twofinite}
Per-fibre finiteness~(b) is combinatorial and unconditional; global
finiteness~(a) is metric and rests on the genuine Hensley/Sylvester degree
bound (Theorem~\ref{thm:hensley}, Theorem~\ref{thm:maxdeg5}). The clean
statement of Q1 is~\eqref{eq:fibrebound}: \emph{the fibre is bounded by the
number of distinct IP-simplex circuits inside one normal-form polytope.}
For the \emph{minimal} reflexive simplices (the duals $\nabla=\Delta^*$,
Definition~\ref{def:min}) the bijection~\eqref{eq:simplexws-bij} makes the
fibre essentially a singleton; the large fibres of $\Phi$ come from
\emph{many minimal simplices feeding one maximal polytope}
(Section~\ref{sec:fibre-loss}).
\end{remark}

\subsection{Q2 --- where the data is lost}\label{sec:fibre-loss}

\begin{proposition}[Both arrows of~\eqref{eq:phifactor} are lossy, for
different reasons]\label{prop:lossy}
In the factorization $q\xrightarrow{\alpha}\Delta_q\xrightarrow{\beta}
\mathrm{NF}$:
\begin{enumerate}
  \item $\alpha$ is injective on the simplicial core
        (\eqref{eq:simplexws-bij}) but \emph{not} on the Newton polytope:
        $\Delta_q$ is the \emph{convex hull} of all degree-$\mathfrak d$
        monomials, and the hull is a lossy non-linear operation --- it discards
        every non-vertex lattice point and can \emph{merge} the contributions
        of different weight systems, so distinct $q$ can yield
        $GL$-equal $\Delta_q$ (the toric analogue of ``different polynomials,
        same Newton polytope'');
  \item $\beta$ quotients by the \emph{non-local} group $GL(d,\mathbb Z)$:
        whether two lattice polytopes are equivalent is not a function of any
        local/arithmetic feature of the weights --- the normal form of
        Section~\ref{sec:algo-nf} is precisely a \emph{certificate} of
        equivalence, and there is no weight-local shortcut to it.
\end{enumerate}
\end{proposition}

\begin{remark}[The deep reason: many minimal seeds, one maximal polytope]
\label{rmk:manyseeds}
The structural source of the collapse is the KS flow
$\nabla\ (\text{minimal})\to\Delta=\nabla^*\ (\text{maximal})\to
\{\text{subpolytopes}\}$ of Sections~\ref{sec:sub}--\ref{sec:dual}. A
\emph{single} maximal reflexive polytope $\Delta$ typically contains
\emph{several inequivalent} IP sub-simplices, i.e.\ several distinct weight
systems, all of which flow to the same $\Delta$ and hence the same
$\mathrm{NF}$ (and the same Calabi--Yau). This is exactly why insertion into the
\texttt{NF\_List} is idempotent (Remark~\ref{rmk:class-code}). Predicting the
fibre in advance therefore means \emph{predicting that two IP circuits of a
polytope one has not yet built lie in the same $GL(d,\mathbb Z)$-orbit} --- which
is the very decision the normal form exists to make. This is the precise sense
in which Q2's ``you cannot tell without computing the polytope'' is correct.
\end{remark}

\begin{remark}[Quantitative loss in $d=5$, reconciled]\label{rmk:fibrequant}
With the corrected counts of Remark~\ref{rmk:datasetnumbers}, the
$\sim$$1.85\times 10^{11}$ reflexive single weight systems collapse to
$\approx 2.3\times 10^{9}$ distinct normal forms, so the \emph{average} fibre of
$\Phi$ over the \emph{full} Sch\"oller--Skarke list is
\begin{equation}\label{eq:avgfibre}
  \overline{\#\Phi^{-1}}\ \approx\ \frac{1.85\times 10^{11}}{2.3\times 10^{9}}
  \ \approx\ 80 ,
\end{equation}
not the $\sim$$2100$ of earlier drafts. The earlier figure divided by the
\emph{sample} normal-form count $8.8\times 10^{7}$ rather than the global
$2.3\times 10^{9}$, and so overstated the collapse by $\sim$$26\times$; it is a
statistic of the sample, not of the classification. \emph{Known vs.\
extrapolated}: the numerator~\eqref{eq:ssnumbers} is a hard Sch\"oller--Skarke
count; the denominator $2.3\times 10^{9}$ is the reported global normal-form
total, of which $8.8\times 10^{7}$ is the portion with a fully materialized
$\mathrm{WS}\to\mathrm{NF}$ map. The honest reading of~\eqref{eq:avgfibre} is
``order $10^{1}$--$10^{2}$ weight systems per polytope on average,'' with the
distribution heavily skewed (most polytopes have small fibres, a few maximal /
high-symmetry ones have very large ones). That spread --- not the mean --- is
the lost data, invisible until one runs $\alpha$ then $\beta$.
\end{remark}

\subsection{Q3 --- non-trivial predictors of the fibre from WS data alone}
\label{sec:fibre-predict}

No weight-local computation decides the \emph{general} fibre
(Proposition~\ref{prop:lossy}\,(2)), but several genuine partial predictors
exist, ordered cheap-to-deep.

\begin{enumerate}
  \item \textbf{Necessary invariants --- but one must split ``WS-local'' from
        ``polytope-level'' (a correction).} It is tempting, and was asserted in
        an earlier draft, that two WS in one fibre must agree on the degree
        $\mathfrak d$ and the weight multiset $\{w_i\}$. \emph{This is false},
        and the materialized $d=5$ map disproves it directly
        (Remark~\ref{rmk:fibredata}): degree and multiset are invariants of the
        \emph{simplicial core} $\nabla_q$, \emph{not} of the Newton-polytope
        normal form, because the hull $\alpha$ merges weight systems of
        \emph{different} degree and \emph{different} weights. The quantities that
        \emph{are} provably constant on a fibre fall into two tiers:
        \begin{itemize}
          \item \emph{WS-local (hull-free), provably fibre-constant}: the count
                of distinct weights $\#\{w_i\}$, the multiplicity profile of the
                multiset, whether any $w_i=1$ and how many, and the Gorenstein
                flag $w_i\mid\mathfrak d$ (Theorem~\ref{thm:reflexweight}).
                These cost nothing but are very coarse (Section~\ref{sec:fibre-sieve}).
          \item \emph{Polytope-level (require $\alpha$)}: the monomial count
                $\ell(\Delta_q)$, the Batyrev Hodge numbers
                $h^{1,1},h^{1,d-1}$, the Euler characteristic, and the
                signature $M{:}np\,nv\ N{:}np\,nv$ (Remark~\ref{rmk:dual-code}).
                These are strong (almost injective on normal forms,
                Remark~\ref{rmk:fibredata}) but are \emph{downstream} of the
                lossy hull, so they are not available in pure weight space.
        \end{itemize}
        \emph{No finite list of either tier is sufficient}: inequivalent
        reflexive polytopes with identical Hodge data and point counts exist.
  \item \textbf{Reflexive-simplex sub-class: fully predictable (Conrads).}
        By~\eqref{eq:simplexws-bij}, $q$ determines $\nabla_q$ up to $GL$, and
        simplex-equivalence is decided by a \emph{much cheaper} normal form than
        the full $\Delta_q$ NF --- essentially the Hermite/Smith form of the
        single circuit relation together with the finite symmetry of the weight
        multiset. On the Gorenstein sub-class $w_i\mid\mathfrak d$ (the
        unit-fraction weight systems, $\sim$$3462$ of them for $d=5$), this is
        Conrads' classification of reflexive simplices by weight systems modulo
        an explicit finite equivalence: \emph{here the fibre is computable from
        the weights with no hull at all.}
  \item \textbf{Lattice-quotient ($\mathbb Z_k$) coincidences.} Two WS defining
        the same toric variety after a finite abelian quotient --- the
        \texttt{z[][]}/\texttt{m[]}/\texttt{nz}/\texttt{index} data of the
        \texttt{CWS} struct (Remark~\ref{rmk:cwsstruct}) and the
        ``$+\mathbb Z_k$'' refinements of Theorem~\ref{thm:cws5types} ---
        coincide for an \emph{arithmetically detectable} reason: compare the
        sublattice index and quotient phases, no hull required.
  \item \textbf{Secondary affine symmetries (oriented-matroid data).}
        Non-permutation symmetries --- affine dependences among the monomial
        generators beyond the defining circuit --- force extra automorphisms of
        $\Delta_q$ (exactly what \texttt{Make\_Poly\_Sym\_NF},
        Remark~\ref{rmk:nf-code}, must quotient) and predict coincidences. These
        live in the sign pattern of the secondary circuits of $q$. Partial and,
        in $d=5$, not fully systematized --- a genuine open thread.
\end{enumerate}

\begin{remark}[The honest negative, confirmed on the data]\label{rmk:fibrenegative}
What \emph{cannot} be done is to decide \emph{general} fibre-membership --- two
IP weight systems with inequivalent simplices but Newton polytopes coinciding
after hulling --- without constructing the polytopes, because that coincidence
is mediated by $\alpha$, which is non-local in the weights
(Proposition~\ref{prop:lossy}\,(1)). The materialized map quantifies just how
non-local: see Remark~\ref{rmk:fibredata}. The reflexive-simplex sub-class
(predictor~2) is the one genuine exception, and is exactly solvable
(Section~\ref{sec:fibre-dichotomy}); the general fibre is a
$GL(d,\mathbb Z)$-orbit decision and is not reducible to a weight-local
computation.
\end{remark}

\begin{remark}[What the materialized $d=5$ map shows]\label{rmk:fibredata}
The \texttt{refpoly-5d} sample (a current snapshot of the
$\mathrm{WS}\to\mathrm{NF}$ map over $\approx 5\times 10^{5}$ distinct normal
forms and $\approx 9.3\times 10^{6}$ reflexive single weight systems) makes the
preceding structure measurable. Three facts:
\begin{enumerate}
  \item[\textnormal{(i)}] \emph{The fibre is heavily skewed, not concentrated at
        its mean.} About $36\%$ of normal forms are \emph{singletons}
        ($\#\Phi^{-1}=1$); the median fibre is $2$; yet the largest observed
        fibre exceeds $6.5\times 10^{5}$ weight systems. Only $\sim$$2\%$ of
        weight systems sit in singleton fibres. The mean of
        Remark~\ref{rmk:fibrequant} is therefore a poor summary; \emph{the
        spread is the phenomenon.}
  \item[\textnormal{(ii)}] \emph{Degree and multiset are \underline{not}
        fibre-constant.} Across multi-WS fibres, the originating degree
        $\mathfrak d$ is constant in only $\sim$$0.15\%$ of fibres and the
        weight multiset in $0\%$ --- direct evidence that $\alpha$ fuses
        weight systems of different degree and shape (the toric ``different
        polynomials, same normal form''). This is the empirical refutation of
        the naive sieve and the reason predictor~1 had to be split.
  \item[\textnormal{(iii)}] \emph{Polytope invariants are nearly injective; the
        cheap WS-local ones are nearly useless.} The full signature
        $(h^{1,1},h^{1,2},h^{1,3},h^{2,2},\chi,\#\text{verts},\#\text{facets},
        \ell,\ell^{*},M{:}np\,nv,N{:}np\,nv)$ separates all but $\sim$$7$ of
        $5\times 10^{5}$ normal forms --- but every component requires $\alpha$.
        The provably fibre-constant \emph{WS-local} invariants of predictor~1,
        used as a sieve, resolve only $\sim$$0.17\%$ of weight-system pairs as
        ``different,'' leaving $99.8\%$ to the hull$+$NF.
\end{enumerate}
The lesson for the stay-in-weight-space goal is stark: \emph{the discriminating
information lives downstream of $\alpha$}, so a complete same-NF test cannot
remain in weight space (Remark~\ref{rmk:fibredesign}), and the genuine
weight-local handle is confined to the simplex sub-class
(Section~\ref{sec:fibre-dichotomy}).
\end{remark}

\subsection{Q4 --- naive checks that two WS give the same polytope}
\label{sec:fibre-naive}

Ordered by cost; the predictors of Section~\ref{sec:fibre-predict} are precisely
attempts to avoid the expensive complete checks here.

\begin{enumerate}
  \item \textbf{Construct $+$ normal-form comparison (complete, canonical).}
        For each $q$ run $\alpha$ (\texttt{Make\_CWS\_Points},
        Remark~\ref{rmk:cwspoints}), reflexivity (\texttt{IP\_Check} /
        \texttt{Ref\_Check}), then \texttt{Make\_Poly\_NF}
        (Remark~\ref{rmk:nf-code}). \emph{Same polytope $\iff$ \texttt{pNF}
        matrices bit-for-bit equal.} This is the definition of the answer and
        what the $88$M dedup does; it pays the full construction $+$ Hermite $+$
        symmetry cost per WS.
  \item \textbf{Construct $+$ brute-force $GL(d,\mathbb Z)$ search.} Build both
        $\Delta_{q_i}$ and search for unimodular $M$ matching vertices. The
        naivest \emph{correct} test, strictly worse than~1: the normal form
        exists precisely to replace this exponential vertex matching by a
        canonical-form comparison.
  \item \textbf{Construct $+$ cheap-invariant prefilter, then NF.} Build both,
        compare $M{:}np\,nv\ N{:}np\,nv$ and Hodge numbers first, run the NF
        only on survivors. The standard practical baseline; the construct-first
        cousin of predictor~1.
  \item \textbf{Oriented vertex/point-set equality (incomplete).} Put both in a
        deterministic but non-canonical orientation (e.g.\ lexicographically
        sorted vertices) and compare. Catches \emph{identical} polytopes but
        \emph{misses} $GL(d,\mathbb Z)$-equivalent-but-not-identical ones, so it
        over-counts the fibre --- the failure mode that makes a true normal
        form~1 necessary.
\end{enumerate}

\begin{remark}[The pipeline design principle]\label{rmk:fibredesign}
Every \emph{complete} naive method (1--2) requires running $\alpha$ (build the
polytope); by Remark~\ref{rmk:fibrenegative} no complete method stays in
weight-space. The invariant methods (3, and predictors~1,3,4 of
Section~\ref{sec:fibre-predict}) are incomplete but cheap. The art of
\texttt{refpoly-5d} is to push maximal discrimination into the cheap
weight-local/invariant sieve so that the expensive complete check~1 runs on as
few within-bucket pairs as possible.
\end{remark}

%--------------------------------------------------------------
\subsection{Why the reflexive-simplex sub-class is fully predictable and
the rest is not}\label{sec:fibre-dichotomy}
%--------------------------------------------------------------

This subsection answers the sharp question directly: \emph{why is the
reflexive-simplex sub-class fully predictable from the weight system alone,
while the general case is not?} The dichotomy is exact and has a one-line cause:
\emph{a simplex has no non-trivial subpolytope combinatorics --- its vertices
\underline{are} the weight data --- whereas a non-simplicial Newton polytope is
the convex hull of many monomials, and its $GL$-class depends on how those
monomials are arranged, which is not a function of the weight multiset.}

\begin{theorem}[Simplex/non-simplex dichotomy for the fibre]
\label{thm:dichotomy}
Let $q$ be a reflexive single weight system in dimension $d$.
\begin{enumerate}
  \item[\textnormal{(S)}] \emph{Simplex case.} If $\Delta_q$ is a $d$-simplex
        (equivalently $\Delta_q=\nabla_q$, i.e.\ the only degree-$\mathfrak d$
        monomials that survive to vertices are the $d{+}1$ fan generators), then
        the $GL(d,\mathbb Z)$-class of $\Delta_q$ is a \emph{closed-form
        function of $q$}: by Lemma~\ref{lem:ipcrit} and the
        bijection~\eqref{eq:simplexws-bij}, $\Delta_q$ is determined up to $GL$
        by its single circuit $\sum_i w_i v_i=0$ together with the index of the
        sublattice the $v_i$ span. Two such $q,q'$ give $GL$-equal $\Delta$ iff
        their circuit data agree up to a permutation and an explicit finite
        lattice automorphism --- this is \textbf{Conrads' classification of
        reflexive simplices by weight systems}. \emph{The fibre is computable
        from the weights with no hull.}
  \item[\textnormal{(N)}] \emph{Non-simplex case.} If $\Delta_q$ has more than
        $d{+}1$ vertices, the extra vertices are degree-$\mathfrak d$ monomials
        that are \emph{not} fan generators, and the face lattice of
        $\operatorname{conv}$ of all monomials is \emph{not} recoverable from the
        weight multiset: the same multiset of weights can produce different
        hull combinatorics, and different multisets can produce the same hull
        (Remark~\ref{rmk:fibredata}\,(ii)). The $GL$-class is then a genuine
        orbit invariant of the realized point configuration, computed only by
        $\beta$ (the normal form), with no weight-local closed form.
\end{enumerate}
\end{theorem}

\begin{proof}[The precise reason]
In case~(S) the polytope has exactly $d{+}1$ vertices and $d{+}1$ facets; there
is a \emph{unique} affine dependence among the vertices
(Lemma~\ref{lem:ipcrit}), and its cleared-denominator coefficient vector
\emph{is} the weight system. Hence ``which vertices, in what lattice'' carries
no information beyond $q$ and the spanning index: a simplex has no proper
sub-circuits, no choice of triangulation, no secondary fan --- nothing for $GL$
to act on except the rigid frame the weights already pin down. The map
$q\mapsto[\Delta_q]_{GL}$ is therefore injective up to the finite Conrads
equivalence, and \emph{invertible in closed form}.

In case~(N) the defining data is the \emph{set} of $\ell(\Delta_q)$ monomials
on the weight hyperplane, and $\Delta_q=\operatorname{conv}$ of that set. The
hull operation forgets every interior and edge-interior monomial and retains
only the extreme ones; \emph{which} monomials are extreme depends on the metric
arrangement of the whole degree-$\mathfrak d$ lattice slice, not on the weight
multiset. Concretely, the extra information the non-simplex case needs and the
simplex case does not is exactly the \emph{oriented matroid of the monomial
point configuration} --- the full set of sign patterns of affine dependences
among the generators (the ``secondary circuits''), which determines the face
lattice and hence the $GL$-orbit. For a simplex this oriented matroid is the
unique uniform one on $d{+}1$ points and contains no information beyond the
circuit; for a hull of $>d{+}1$ monomials it is a rich, weight-non-local object.
This is the precise content of predictor~4 of
Section~\ref{sec:fibre-predict}. $\square$
\end{proof}

\begin{remark}[Restating the dichotomy in one sentence]\label{rmk:dichotomy}
\emph{In the simplex case the vertices are the weight data; in the non-simplex
case the vertices are a hull of many monomials whose combinatorics is extra,
weight-non-local information.} That single difference is the whole reason
predictor~2 is exact and predictors~1,3,4 are only partial. The Gorenstein
unit-fraction weight systems ($w_i\mid\mathfrak d$, $\sim$$3462$ of them in
$d=5$, OEIS A002966) are the cleanest inhabitants of case~(S); but case~(S) is
\emph{not} the same as ``Gorenstein'' --- it is the geometric condition
``$\Delta_q$ is a simplex,'' which one can read off cheaply by checking that
$\Delta_q$ has exactly $d{+}1$ vertices (a single hull, far cheaper than the
full NF).
\end{remark}

%--------------------------------------------------------------
\subsection{Generalizing Conrads by Gale duality: the right objects, the
genuine obstruction, and the ceiling}\label{sec:fibre-gale}
%--------------------------------------------------------------

Theorem~\ref{thm:dichotomy} draws the line where weight-space reasoning stops.
This subsection answers the two sharp research questions it raises.
\emph{Q1: what objects would \underline{generalize} Conrads' simplex method to
the non-simplicial case?} \emph{Q2: what is the \underline{best} one could hope
for --- a non-trivial relation deciding same-fibre membership without the normal
form, ideally without even the convex hull of $\Delta_q$?} We frame both through
\textbf{Gale duality}, isolate the one step that provably cannot stay in weight
space, and stratify the answer by the \emph{Gale dimension}. We close with an
empirical calibration on the full $2.3\times10^{9}$-polytope list that tells the
reader exactly how much of \emph{their} data the tractable strata reach.

\subsubsection*{Q1 --- Gale duality is the correct generalization of ``the
weight system of a simplex''}

The starting observation makes Conrads' method look like the bottom rung of a
ladder rather than an isolated trick.

\begin{proposition}[A (combined) weight system presents the space of affine
relations]\label{prop:wsgale}
Let $A=\{a_1,\dots,a_n\}\subset\mathbb Z^{d}$ be the vertex configuration of a
polytope in $N_{\mathbb R}\cong\mathbb R^d$ (for us $A=\operatorname{Vert}
\Delta_q$, sitting on the degree-$\mathfrak d$ slice), written projectively as
the columns of
\begin{equation}\label{eq:Ahat}
  \widehat A=\begin{pmatrix} 1 & \cdots & 1\\ a_1 & \cdots & a_n\end{pmatrix}
  \in\mathbb Z^{(d+1)\times n}.
\end{equation}
The \emph{space of affine dependences} (the linear relations
$\sum_i\lambda_i a_i=0$, $\sum_i\lambda_i=0$) is $\ker\widehat A$, of dimension
\begin{equation}\label{eq:galedim}
  g \;=\; n-d-1 \;=:\; \text{the \emph{Gale dimension}}.
\end{equation}
A choice of basis of $\ker\widehat A$, recorded as the rows of a matrix
$B\in\mathbb Z^{g\times n}$ with $\widehat A B^{\mathsf T}=0$, is a
\emph{Gale transform} (Gale dual) $\widehat A^{\vee}=\{b^1,\dots,b^n\}\subset
\mathbb R^{g}$ of $A$. A single weight system $w=(w_0,\dots,w_d)$ is exactly the
case $n=d+1$, $g=1$: the one-dimensional relation space is spanned by the single
circuit $\sum_i w_i v_i=0$ (Lemma~\ref{lem:ipcrit}), and ``the weight system''
\emph{is} that $1$-row $B=(w_0,\dots,w_d)$. A combined weight system is a
spanning set / partial basis of (part of) the same relation space: its $k$ rows
are $k$ independent circuits, so it presents a $k$-dimensional subspace of
$\ker\widehat A$.
\end{proposition}

\begin{proof}
$\widehat A$ has full row rank $d+1$ (the $a_i$ affinely span), so $\dim\ker
\widehat A=n-(d+1)=g$. The first row of $\widehat A$ forces $\sum_i\lambda_i=0$,
i.e.\ $\ker\widehat A$ is precisely the affine (not merely linear) dependences.
For $n=d+1$ this kernel is a line; clearing denominators of any nonzero vector in
it gives the coprime circuit, whose entries are the $w_i$ of
Lemma~\ref{lem:ipcrit}. The CWS statement is Definition~\ref{def:combined}
re-read: each block circuit $\sum_i w^{(j)}_i v_i=0$ is an element of
$\ker\widehat A$, and complementary-span (Theorem~\ref{thm:cwsvalid}) is linear
independence of the corresponding rows. $\square$
\end{proof}

This reframes the whole problem. ``The weight system is an invariant of the
simplex, not extra data'' (Remark~\ref{rmk:twofinite}) is the $g=1$ instance of
the general fact that \emph{the Gale transform is a $GL(d,\mathbb Z)$-invariant
of the configuration}: $A\mapsto MA$ for $M\in GL(d,\mathbb Z)$ leaves
$\ker\widehat A$ unchanged, so the Gale dual is unchanged up to a change of basis
$B\mapsto UB$, $U\in GL(g,\mathbb Z)$. The natural object generalizing ``the
weight system of a simplex'' is therefore the Gale transform of the vertex
configuration, \emph{decorated} with the two pieces of structure Conrads uses
implicitly at $g=1$:

\begin{definition}[The decorated Gale datum]\label{def:galedatum}
For a vertex configuration $A$ with Gale dual $B\in\mathbb Z^{g\times n}$, the
\emph{decorated Gale datum} is the triple
\begin{equation}\label{eq:galedatum}
  \bigl(\ \underline{\chi_B}\ ,\ \underline{\mathrm{SNF}}\ ,\ H\ \bigr),
\end{equation}
\begin{enumerate}
  \item[\textnormal{(a)}] the \emph{oriented matroid} (chirotope) $\chi_B$ of the
        Gale dual --- the list of signs of all maximal minors of $B$, equivalently
        the sign vectors of all relations $\sum_i\lambda_i a_i=0$. This is the
        affine-over-$\mathbb R$ combinatorial type of $A$: by the Gale criterion a
        subset $F\subset A$ is a \emph{face} of $\operatorname{conv}A$ iff its
        complement $A\setminus F$ is a \emph{positive (co)circuit} of the Gale
        dual (the complementary points positively span $\mathbb R^g$). So
        $\chi_B$ determines the entire face lattice;
  \item[\textnormal{(b)}] the \emph{arithmetic decoration}: the Smith normal
        forms of the circuit matrices (equivalently the elementary divisors of
        $B$ and of the sublattice the relations cut out), recording the local
        indices $[\,N : \langle\text{relevant sublattice}\rangle\,]$;
  \item[\textnormal{(c)}] the \emph{torsion subgroup}
        $H=N/\langle A\rangle_{\mathbb Z}$ (the finite quotient of the ambient
        lattice by the sublattice generated by the configuration) --- the
        ``$+\mathbb Z_k$'' data of the \texttt{CWS} struct
        (Remark~\ref{rmk:cwsstruct}).
\end{enumerate}
\end{definition}

\begin{theorem}[What the non-simplex case needs that the simplex case does not]
\label{thm:galerefine}
Equivalence of vertex configurations refines in two stages:
\begin{equation}\label{eq:refinechain}
  GL(d,\mathbb Z)\text{-equivalence}
  \ \subsetneq\
  \text{oriented-matroid (chirotope) equivalence},
\end{equation}
read as: $GL(d,\mathbb Z)$-equivalent configurations have the same chirotope, but
the converse \emph{fails}. The gap is exactly the arithmetic decoration
(b)--(c): two configurations with the same chirotope $\chi_B$ can have different
lattice indices/torsion and hence be $GL(d,\mathbb Z)$-inequivalent. At $g=1$
(the simplex) the chirotope is the unique uniform oriented matroid on $d{+}1$
points and carries \emph{no} information beyond ``it is a simplex''; there the
entire content is arithmetic, and the decoration collapses to the single Smith
form of the circuit $(w_0,\dots,w_d)$ --- which is precisely
\textbf{Conrads' invariant}. For $g\ge1$ with $>d{+}1$ points the chirotope is
genuine combinatorial data \emph{and} the arithmetic decoration is needed on top
of it.
\end{theorem}

\begin{proof}[Mechanism]
A unimodular $M\in GL(d,\mathbb Z)$ permutes nothing and preserves all signed
affine dependences, hence preserves $\chi_B$; this is the inclusion. That it is
strict is exactly the false-positive phenomenon already exhibited in
Example~\ref{ex:fp1}: there $q=(1,1,1,1,1,1)$ and $q'=(1,1,2,2,2,4)$ are both
$5$-simplices --- \emph{identical} chirotope (the uniform one on $6$ points),
identical face lattice, identical normalized volume --- yet $GL(5,\mathbb Z)$-
inequivalent, separated only by the circuit Smith data $(1,1,1,1,1,1)$ vs.\
$(1,1,2,2,2,4)$. That separating datum is the arithmetic decoration (b). The
$g=1$ collapse is Lemma~\ref{lem:ipcrit}: one circuit, one Smith form, no
combinatorial freedom. $\square$
\end{proof}

\begin{remark}[The generalization, stated as a program]\label{rmk:galeprogram}
Conrads at $g=1$ inputs the circuit and outputs a closed-form $GL$-class because
\emph{the chirotope is trivial and only the arithmetic remains}. The honest
generalization to $g\ge1$ is: \emph{(i)} compute the chirotope $\chi_B$ of the
Gale dual (this fixes the face lattice and the $GL$-orbit \emph{up to} the
arithmetic refinement); \emph{(ii)} compute the Smith/torsion decoration
(b)--(c) to resolve the residual arithmetic ambiguity within a chirotope class.
The pair $(\chi_B,\text{decoration})$ is a \emph{complete} $GL(d,\mathbb Z)$-
invariant of the configuration. Nothing here is conjectural --- it is a faithful
restatement of the orbit problem. The catch is entirely in the \emph{input}:
one must already know $A=\operatorname{Vert}\Delta_q$. That is the obstruction.
\end{remark}

\subsubsection*{Q2, part 1 --- the genuine obstruction: which step needs the hull}

\begin{theorem}[Weight-local vs.\ hull-requiring: a clean separation]
\label{thm:obstruction}
Run the program of Remark~\ref{rmk:galeprogram} on a single weight system $q$.
Then:
\begin{enumerate}
  \item[\textnormal{(W)}] \emph{Genuinely weight-local} (no hull; only linear
        algebra over $\mathbb Q$ on the relation matrix and integer
        Smith-normal-form / index computations): everything \emph{downstream of
        knowing the vertex set}. Given $A=\operatorname{Vert}\Delta_q$ as input,
        the Gale dual $B$, its chirotope $\chi_B$, the Smith decoration, and the
        torsion $H$ are all polynomial-time linear-algebraic functions of $A$ ---
        no convexity, no normal form.
  \item[\textnormal{(H)}] \emph{Provably requires the hull}: the determination of
        \emph{which} degree-$\mathfrak d$ lattice points of $q$ are
        \underline{vertices} of $\Delta_q$. Selecting $\operatorname{Vert}
        \Delta_q$ from the $\ell(\Delta_q)$ monomials is exactly the
        \emph{convex-hull / per-point LP-feasibility} problem ($x$ is a vertex iff
        $x\notin\operatorname{conv}(\text{other points})$, a linear program), and
        it is not a function of the weight multiset
        (Theorem~\ref{thm:dichotomy}(N), Remark~\ref{rmk:fibredata}(ii)).
\end{enumerate}
Consequently the irreducible residue, after all weight-local work, is the
\textbf{lattice-polytope $GL(d,\mathbb Z)$-isomorphism problem on the
vertex configuration} --- and its only weight-non-local ingredient is the vertex
selection~(H).
\end{theorem}

\begin{proof}[Where it bites]
Statement (W) is Theorem~\ref{thm:galerefine}: the chirotope is the sign data of
minors of $B$, $B$ is a kernel basis of $\widehat A$, and Smith forms and the
torsion $N/\langle A\rangle$ are integer linear algebra --- all polynomial in the
bit-size of $A$. Statement (H) is the content of the lossy arrow $\alpha$
(Proposition~\ref{prop:lossy}(1)): $\Delta_q$ is $\operatorname{conv}$ of all
degree-$\mathfrak d$ monomials, and the hull retains only the extreme ones; which
monomials are extreme depends on the metric arrangement of the whole lattice
slice, not on $(w_0,\dots,w_d)$. There is no closed form for ``vertex or not''
that avoids a feasibility test, and Remark~\ref{rmk:fibredata}(ii) shows
empirically that even $\#\operatorname{Vert}$ is not a function of the weight
multiset. Hence the one step the program cannot import from weight space is
exactly the one Q2 wished to avoid. $\square$
\end{proof}

\begin{remark}[Reading the obstruction honestly]\label{rmk:obstruction}
This is the precise, and slightly deflationary, answer to Q2's strongest form
(``without even the convex hull''): \emph{the hull is unavoidable, because the
vertex set is the input to every invariant downstream, and the vertex set is the
hull.} What \emph{is} salvaged is sharp and useful: once the (single, unavoidable)
hull has produced $\operatorname{Vert}\Delta_q$, the entire $GL$-classification
is weight-local linear/arithmetic algebra via the Gale datum --- one never needs
the expensive symmetry-reduced normal form $\beta$ of Section~\ref{sec:algo-nf}.
This is exactly the gain already cashed out in sieve stage~\textbf{S3}
(Section~\ref{sec:fibre-sieve}) for $g=1$; the program extends that gain to all
strata, replacing $\beta$ by ``(one hull) $+$ Gale datum.''
\end{remark}

\subsubsection*{Q2, part 2 --- the ceiling: a Gale-dimension hierarchy}

The best achievable is therefore a \emph{layered} statement, stratified by the
Gale dimension $g=n-d-1$ of~\eqref{eq:galedim}. Each layer is a known chapter of
the theory of \emph{polytopes with few vertices} (Gale diagrams).

\begin{theorem}[The few-vertex / Gale-dimension ceiling]\label{thm:ceiling}
Let $q$ be a reflexive single weight system in dimension $d$ with
$\Delta_q$ having $n$ vertices, $g=n-d-1$.
\begin{enumerate}
  \item[$g=0$ \textnormal{($n=d{+}1$, simplices):}] Conrads is \emph{complete and
        closed-form}: $\Delta_q$ is determined up to $GL(d,\mathbb Z)$ by the
        single circuit's Smith form (Theorem~\ref{thm:dichotomy}(S)). No hull
        beyond confirming $n=d{+}1$. \emph{Done.}
  \item[$g=1$ \textnormal{($n=d{+}2$):}] The Gale diagram is a configuration of
        $n$ signs $\{+,-,0\}$ on a line --- the classical Gale-diagram description
        of polytopes with $d{+}2$ vertices (Gr\"unbaum). The chirotope is fixed by
        the multiset of signs, so the face lattice is a closed-form function of
        the sign pattern of the single nontrivial relation among the vertices.
        \emph{Claim (the concrete novel target):} the $GL(d,\mathbb Z)$-class is a
        closed-form function of the Gale sign pattern together with the arithmetic
        decoration (b)--(c) of Definition~\ref{def:galedatum} --- i.e.\ Conrads
        \emph{genuinely generalizes} here, with no normal form, only the Smith
        data of the (now two-dimensional) relation lattice.
  \item[$g=2$ \textnormal{($n=d{+}3$):}] The Gale diagram is an \emph{affine}
        point configuration in the plane (Gale's ``standard diagram''); the
        combinatorial types of $d{+}3$-vertex polytopes are still classifiable in
        principle, so the chirotope is enumerable and the program
        (Remark~\ref{rmk:galeprogram}) remains tractable.
  \item[$g\ge3$:] The oriented-matroid classification grows wildly (realizability
        becomes $\exists\mathbb R$-complete in general); no closed form is
        available, and one falls back on a \emph{sound-but-incomplete} invariant
        sieve plus the normal form on the residue.
\end{enumerate}
\end{theorem}

\begin{remark}[The honest ``best we can hope for'']\label{rmk:ceiling}
Assembling Theorem~\ref{thm:ceiling}: the optimal attainable object is
\emph{a polynomial-time, $GL(d,\mathbb Z)$-invariant computable from the (C)WS
relation matrix plus a single hull, that is}
\begin{itemize}
  \item \emph{complete} on the structurally characterized low-Gale-dimension
        strata $g\le 2$ (the generalized Conrads: chirotope of a $\le2$-dim Gale
        diagram $+$ arithmetic decoration), \emph{replacing $\beta$ entirely};
  \item a \emph{refined sound sieve} for $g\ge3$ (the fibre-constant WS-local
        invariants of predictor~1, $+$ the near-injective polytope signature S5),
        which decides ``different'' cheaply but cannot certify ``same'';
  \item with a \emph{provable irreducible residue}: the within-chirotope,
        $g\ge3$ pairs, whose separation is the genuine lattice-polytope
        $GL(d,\mathbb Z)$-isomorphism problem (Theorem~\ref{thm:obstruction}).
\end{itemize}
A clean impossibility statement is as valuable as the positive one: \emph{no
weight-local computation can decide same-fibre membership in the $g\ge3$ residue
without the hull}, because the discriminating data (the vertex set, hence the
chirotope) is created by $\alpha$ and is provably not a function of the weights
(Theorem~\ref{thm:obstruction}, Remark~\ref{rmk:fibredata}). Q2 has a sharp
two-part answer: \emph{yes} below $g=3$ (and even without $\beta$), \emph{no}
above it (not even without $\alpha$).
\end{remark}

\begin{remark}[Surrounding theory to draw on]\label{rmk:surrounding}
The program sits inside well-developed machinery the reader can pull on directly.
\emph{(i)} The \textbf{GKZ secondary polytope / secondary fan} of the point
configuration $A$ (Gelfand--Kapranov--Zelevinsky) organizes \emph{all} regular
triangulations of $A$; its rays are the circuits, so the (C)WS data are
coordinates on the secondary fan and the ``secondary circuits'' of predictor~4
(Section~\ref{sec:fibre-predict}) are literally its $1$-skeleton. \emph{(ii)} The
\textbf{$A$-discriminant} (GKZ) is the locus in coefficient space where $\alpha$'s
hull degenerates --- the analytic shadow of the vertex-selection problem~(H).
\emph{(iii)} The \textbf{Gale diagrams of polytopes with few vertices}
(Gr\"unbaum; Sturmfels) are exactly the $g\le2$ classification of
Theorem~\ref{thm:ceiling}. \emph{(iv)} The \textbf{lattice/arithmetic refinement
of oriented matroids} --- chirotope plus elementary-divisor data --- is the
formal home of the decoration (b)--(c) and of inclusion~\eqref{eq:refinechain}.
\end{remark}

\subsubsection*{Q2, part 3 --- empirical calibration: does this reach the
Sch\"oller--Skarke data?}

The hierarchy is the right \emph{ceiling}; whether it is worth building for
\emph{this} dataset is an empirical question, because Sch\"oller--Skarke
enumerates \emph{maximal} reflexive polytopes (Remark~\ref{rmk:datasetnumbers}),
which --- being duals of \emph{minimal} ones --- tend to have \emph{many}
vertices, i.e.\ \emph{large} Gale dimension. We measured the vertex-count
distribution over the \emph{entire} classification.

\begin{remark}[Vertex-count / Gale-dimension distribution of the full $d=5$
list]\label{rmk:galecalib}
Over all $N_{\mathrm{NF}}=2\,377\,115\,412$ distinct normal forms and
$N_{\mathrm{WS}}=185\,269\,499\,015$ reflexive single weight systems (the
materialized $\mathrm{WS}\to\mathrm{NF}$ table; per-vertex breakdown measured
directly), with $d=5$ so $g=n_v-6$:
\begin{center}\small
\begin{tabular}{r r r r r r}
\hline
$n_v$ & $g$ & \#NF & \% of NF & \#WS (fibre mass) & \% of WS \\
\hline
$6$  & $0$ & $50\,001$        & $0.0021\%$ & $1.31\times10^{9}$ & $0.71\%$ \\
$7$  & $1$ & $378\,704$       & $0.0159\%$ & $4.87\times10^{9}$ & $2.63\%$ \\
$8$  & $2$ & $1\,856\,608$    & $0.0781\%$ & $9.90\times10^{9}$ & $5.35\%$ \\
$9$  & $3$ & $6\,489\,821$    & $0.273\%$  & $1.47\times10^{10}$& $7.92\%$ \\
$12$ & $6$ & $7.30\times10^{7}$& $3.07\%$  & $2.09\times10^{10}$& $11.25\%$ \\
$16$ & $10$& $2.38\times10^{8}$& $10.0\%$  & $1.23\times10^{10}$& $6.65\%$ \\
$17$ & $11$& $2.52\times10^{8}$& $10.6\%$  & $9.45\times10^{9}$ & $5.10\%$ \\
\hline
\end{tabular}
\end{center}
(The full table runs $n_v=6,\dots,46$; the distribution peaks at $n_v=17$ by NF
count and $n_v=12$ by WS mass, with median $n_v\approx17$.) The cumulative
tractable strata are
\begin{center}\small
\begin{tabular}{l r r}
\hline
stratum & \% of NF & \% of WS \\
\hline
$g=0$ ($n_v=6$, simplices, Conrads) & $0.0021\%$ & $0.71\%$ \\
$g\le1$ ($n_v\le7$)                 & $0.018\%$  & $3.34\%$ \\
$g\le2$ ($n_v\le8$)                 & $0.096\%$  & $8.68\%$ \\
$g\le5$ ($n_v\le11$)                & $2.77\%$   & $37.5\%$ \\
$g\le10$ ($n_v\le16$)               & $36.3\%$   & $84.1\%$ \\
\hline
\end{tabular}
\end{center}
\end{remark}

\begin{remark}[What the calibration means for the stay-in-weight-space goal]
\label{rmk:galecalib-read}
Two readings, and they point opposite ways --- this is the honest punchline.
\begin{enumerate}
  \item[\textnormal{(i)}] \emph{By NF count the program is marginal.} The
        closed-form strata $g\le2$ contain only $0.096\%$ of the distinct
        polytopes; the median polytope has $n_v\approx17$ ($g\approx11$), deep in
        the explosive regime where Remark~\ref{rmk:ceiling}'s impossibility bites.
        If the goal is to classify \emph{NFs}, the Gale program resolves a sliver
        and the $GL(5,\mathbb Z)$ normal form remains unavoidable for the bulk.
  \item[\textnormal{(ii)}] \emph{By weight-system mass it is far from marginal.}
        The same $g\le2$ strata carry $8.7\%$ of all weight systems, and
        $g\le5$ carries $37.5\%$, because low-vertex polytopes have enormous
        fibres (mean fibre $\approx26\,000$ at $n_v=6$, $\approx12\,900$ at
        $n_v=7$, versus $\approx 80$ overall). Strikingly, the
        \emph{single largest fibre in the entire classification} ---
        $650\,642\,665$ weight systems collapsing to one polytope, the maximal-
        $h^{1,1}$ ($h^{1,1}=28\,348$) Calabi--Yau fourfold of weight system
        $(62,81,491,1206,1647,1840)$ --- has $n_v=7$, i.e.\ \emph{Gale dimension
        $1$}. Since the dedup cost is paid \emph{per weight-system pair}, a
        generalized-Conrads predictor on $g\le2$ would short-circuit the
        $\beta$-step for nearly a tenth of all weight systems, concentrated in
        exactly the giant fibres that dominate the comparison load.
\end{enumerate}
The verdict: the Gale-stratification program is the correct theoretical ceiling
and is \emph{operationally worthwhile precisely because the fibre mass is
front-loaded onto the few-vertex strata}, even though those strata are a
vanishing fraction of the polytopes. It does \emph{not} dissolve the residue ---
the high-$g$ bulk provably needs the normal form (Remark~\ref{rmk:ceiling}) ---
but it removes the heaviest fibres from the expensive path. (Methodological note:
the per-row \texttt{vertex\_count} column of the $5\times10^{5}$ \emph{sample}
parquet is unreliable; the true $n_v$ here is read from the normal-form vertex
matrix \texttt{nf\_vertices} and cross-checked against the full-classification
per-vertex tally.)
\end{remark}

%--------------------------------------------------------------
\subsection{Predictor 4 worked out: a positive coincidence and two honest
false positives}\label{sec:fibre-worked}
%--------------------------------------------------------------

We make predictor~4 (Section~\ref{sec:fibre-predict}) and the dichotomy
(Theorem~\ref{thm:dichotomy}) concrete on small, hand-checkable $d=5$ weight
systems. All Newton polytopes below were verified by exact integer arithmetic
(facet normals over $\mathbb Q$); ``$\#$verts / $\#$facets'' and the facet
lattice-distances are exact.

\begin{example}[Positive: a coincidence detectable from the WS alone]
\label{ex:positive}
Let $q=(1,1,1,1,2,2)$ and $q'=(1,1,2,2,1,1)$, both of degree $\mathfrak d=8$.
They are the \emph{same weight multiset} in a different coordinate order. The
coordinate permutation $\pi$ carrying $q$ to $q'$ acts on the monomial lattice
$\mathbb Z^{6}$ by permuting coordinates --- a \emph{lattice automorphism} --- and
sends the degree-$8$ monomials of $q$ bijectively onto those of $q'$. Hence
$\Delta_{q}$ and $\Delta_{q'}$ are equal after this unimodular map:
$\mathrm{NF}(\Delta_q)=\mathrm{NF}(\Delta_{q'})$, \emph{provably and from the
weights alone}, with no hull. This is the firing case of predictor~4 in its
simplest guise (the automorphism is a coordinate permutation); the genuine
content of predictor~4 is the \emph{non}-permutation lattice automorphisms,
which require the oriented-matroid data of
Theorem~\ref{thm:dichotomy}(N). Empirically (Remark~\ref{rmk:fibredata}) the
multiplicity profile, here $\{1{:}4,\,2{:}2\}$, is one of the few provably
fibre-constant WS-local invariants --- consistent with permutation-coincidences
being the cleanest detectable collapses.
\end{example}

\begin{example}[False positive I: a strong fingerprint that still fails]
\label{ex:fp1}
Consider the two reflexive single weight systems
\[
  q=(1,1,1,1,1,1)\ (\mathfrak d=6),
  \qquad
  q'=(1,1,2,2,2,4)\ (\mathfrak d=12).
\]
Both Newton polytopes are reflexive $5$-\emph{simplices}: each has
\emph{exactly} $6$ vertices, $6$ facets all at lattice distance $1$, and the
\emph{same} normalized volume $5!\cdot\operatorname{vol}=7776=6^{5}$ --- they
agree on every coarse and even most fine invariants one might hope to use.
\emph{Yet they are not $GL(5,\mathbb Z)$-equivalent.} The separating invariant is
exactly the one Theorem~\ref{thm:dichotomy}(S) singles out: the intrinsic
circuit weight (the unique positive relation among the six vertices and the
interior point) is $(1,1,1,1,1,1)$ for $q$ but $(1,1,2,2,2,4)$ for $q'$. So on
the simplex sub-class the cheap circuit normal form decides correctly where
volume, vertex- and facet-counts all collide --- which is precisely why
predictor~2 is the right tool there, and why a generic fingerprint (vertices,
facets, volume, facet-distances) is \emph{not} enough.
\end{example}

\begin{example}[False positive II: equal cheap WS invariants, different polytope]
\label{ex:fp2}
Consider
\[
  q=(1,1,1,3,4,4),
  \qquad
  q'=(1,1,2,2,2,6),
  \qquad \text{both of degree } \mathfrak d=14 .
\]
They have the \emph{identical} two cheapest polytope-blind invariants: degree
$14$ and exactly $651$ degree-$\mathfrak d$ monomials. A naive sieve keying on
$(\mathfrak d,\ell)$ would bucket them together. But their Newton polytopes
differ: $\Delta_q$ has $(\#\text{verts},\#\text{facets})=(14,9)$ and is
\emph{not even reflexive} (a facet sits at lattice distance $5$), while
$\Delta_{q'}$ has $(10,7)$ and \emph{is} reflexive (all facets at distance $1$).
This is the honest limit of any purely WS-local predictor: matching cheap
invariants is necessary, never sufficient, and here the discrepancy only
surfaces after $\alpha$. (It also re-illustrates the $d=5$ watershed of
Section~\ref{sec:fibre-three}: $\Delta_q$ is IP but non-reflexive, the failure
mode that \texttt{Ref\_Check} exists to catch.)
\end{example}

\begin{remark}[Reading the three examples together]\label{rmk:worked}
Example~\ref{ex:positive} shows predictor~4 \emph{can} certify a coincidence
from the weights alone --- when a lattice automorphism (here a permutation) is
visible in the WS. Example~\ref{ex:fp1} shows that on the simplex sub-class the
\emph{correct} certificate is the circuit weight (predictor~2), not a metric
fingerprint. Example~\ref{ex:fp2} shows that off the simplex sub-class no cheap
WS-local invariant is decisive. Together they delimit exactly where weight-space
reasoning succeeds (case~S, Theorem~\ref{thm:dichotomy}) and where it cannot
(case~N).
\end{remark}

%--------------------------------------------------------------
\subsection{A prioritized sieve for the Sch\"oller--Skarke $d=5$ list}
\label{sec:fibre-sieve}
%--------------------------------------------------------------

We synthesize the above into an ordered pipeline that buckets weight systems
into \emph{provably-same}, \emph{provably-different}, and
\emph{undecided-without-NF}, maximizing the fraction decided \emph{before}
paying for $\alpha+\beta$. The numerical resolution figures are measured on the
materialized map (Remark~\ref{rmk:fibredata}).

\begin{enumerate}
  \item[\textbf{S0}] \textbf{Canonicalize the WS (free).} Sort the weights
        ascending and reduce by $\gcd$ (already $1$ on the Sch\"oller--Skarke
        list). \emph{Output:} a canonical weight tuple; two WS with the same
        canonical tuple are \emph{identical}, not merely fibre-mates.
  \item[\textbf{S1}] \textbf{Provably-same by permutation (free, \emph{decides
        same}).} Group WS by their sorted weight \emph{multiset}. Within a
        group, any two are related by a coordinate permutation, hence give the
        same $\mathrm{NF}$ (Example~\ref{ex:positive}). \emph{This is the only
        cheap test that \underline{certifies} same-NF.} Caution: it does
        \emph{not} find all same-NF pairs, because the multiset is \emph{not}
        fibre-constant (Remark~\ref{rmk:fibredata}(ii)); different multisets can
        still share an NF.
  \item[\textbf{S2}] \textbf{Provably-different by fibre-constant WS-local
        invariants (free, \emph{decides different}).} Compute the vector of
        \emph{provably fibre-constant} WS-local invariants:
        $\#\{w_i\}$ (distinct weights), the multiplicity profile, $\#\{i:w_i=1\}$,
        and the Gorenstein flag $w_i\mid\mathfrak d$. Two WS that differ on this
        vector are in \emph{different} fibres. \emph{Honest yield:} on the
        materialized map this resolves only $\sim$$0.17\%$ of pairs --- coarse,
        but free and sound. (Degree, $\min$, $\max$, $\#$Gorenstein slots are
        \emph{not} fibre-constant and must \underline{not} be used here.)
  \item[\textbf{S3}] \textbf{Simplex sub-class: solve exactly (one hull, no NF).}
        Build $\Delta_q$ once and test whether it is a $d$-simplex
        ($\#\text{verts}=d{+}1$); if so, compute the intrinsic circuit weight
        and apply \textbf{Conrads' equivalence} (Theorem~\ref{thm:dichotomy}(S),
        Example~\ref{ex:fp1}). \emph{This decides the fibre completely for the
        simplex sub-class} at the cost of a single hull and an integer normal
        form of one $(d{+}1)$-vector --- far cheaper than the full
        $\mathrm{NF}(\Delta_q)$, and it correctly separates the volume-colliding
        pairs that defeat metric fingerprints.
  \item[\textbf{S4}] \textbf{$\mathbb Z_k$ lattice-quotient coincidences
        (arithmetic).} For WS carrying a sublattice quotient (the
        $\texttt{z}/\texttt{m}/\texttt{nz}/\texttt{index}$ data of the
        \texttt{CWS} struct, Remark~\ref{rmk:cwsstruct}), compare the index and
        quotient phases; equal quotient data on otherwise-matching circuits
        forces the same NF without a further hull (predictor~3).
  \item[\textbf{S5}] \textbf{Polytope-invariant prefilter (one hull each,
        \emph{decides different}).} For the non-simplex residue, build $\Delta_q$
        and bucket by the strong signature
        $(\ell,\ell^{*},\#\text{verts},\#\text{facets},h^{1,1},h^{1,2},h^{1,3},
        h^{2,2},\chi,M{:}np\,nv,N{:}np\,nv)$. Different signatures $\Rightarrow$
        different NF; this signature is \emph{nearly injective}
        (Remark~\ref{rmk:fibredata}(iii)), so the within-bucket residue is tiny.
        These invariants are \emph{not} WS-local --- they cost one $\alpha$ ---
        but no $\beta$.
  \item[\textbf{S6}] \textbf{Normal form (the certificate, \emph{decides
        both}).} On the within-signature residue only, compute
        $\mathrm{NF}(\Delta_q)$ (Section~\ref{sec:algo-nf}) and compare
        bit-for-bit. This is the irreducible step: by
        Proposition~\ref{prop:lossy}(2) and Remark~\ref{rmk:fibredata} the
        general fibre is a $GL(d,\mathbb Z)$-orbit decision with no weight-local
        shortcut.
\end{enumerate}

\begin{remark}[What each stage can and cannot resolve]\label{rmk:sieveyield}
The sieve cleanly separates the sound cheap tests from the expensive complete
one:
\begin{itemize}
  \item \textbf{Decided in pure weight space (S0--S2,S4):} exact \emph{equality}
        of WS (S0); \emph{same}-NF by permutation (S1); \emph{different}-NF by
        fibre-constant invariants (S2, $\sim$$0.17\%$ of pairs) and by quotient
        data (S4). Sound but \emph{incomplete} --- by
        Remark~\ref{rmk:fibredata} the vast majority of pairs survive to a hull.
  \item \textbf{Decided with one hull, no NF (S3,S5):} the entire \emph{simplex
        sub-class} (S3, exactly, via Conrads) and \emph{most} of the rest as
        \emph{different} via the near-injective signature (S5).
  \item \textbf{Irreducible residue (S6):} the within-signature, non-simplex
        pairs --- genuinely different weight systems whose Newton polytopes
        coincide after hulling. \emph{This is the only part that truly needs the
        $GL(d,\mathbb Z)$ normal form}, and it is irreducibly so.
\end{itemize}
The design goal --- ``maximize the fraction decided without the NF'' --- is met
by S1--S5; the honest conclusion is that a \emph{complete} same-NF decision
cannot stay in weight space, and the cheapest complete sieve is
``permutation $+$ one hull $+$ near-injective signature $+$ NF on the small
residue.''
\end{remark}

%--------------------------------------------------------------
\subsection{Computational realization and verification: the Gale fibre key on
the materialized $88.6\text{M}$-fibre $d=5$ list}\label{sec:fibre-verify}
%--------------------------------------------------------------

We implemented the program of Section~\ref{sec:fibre-gale} as a concrete,
hashable \emph{Gale fibre key} and tested it against ground truth on the
materialized $\mathrm{WS}\to\mathrm{NF}$ sample: a table with \emph{one row per
distinct normal form}, $N=88{,}588{,}676$ fibres in total (a sample of the full
$2.377\times10^{9}$; each row carries its complete weight-system fibre and the
PALP normal-form vertex matrix \texttt{nf\_vertices}). The experiment answers, at
scale and against the true NF, the operational question behind
Remark~\ref{rmk:ceiling}: \emph{how complete is the ``(one hull) $+$ Gale
datum'' replacement for $\beta$, and where exactly does it stop?}

\subsubsection*{The key as a layered $GL(5,\mathbb Z)$-invariant}

Given $A=\operatorname{Vert}\Delta_q$ (the single, unavoidable hull of
Theorem~\ref{thm:obstruction}), the key is the cumulative tuple
\begin{equation}\label{eq:keylayers}
  \underbrace{(g,n_v,n_F,\operatorname{vol})}_{\mathrm{L0}}\ \prec\
  \underbrace{+\,\text{facet-size \& vertex-degree multisets}}_{\mathrm{L1}\ (\text{combinatorial type})}\ \prec\
  \underbrace{+\,\text{vertex--facet pairing multisets}}_{\mathrm{Lpm}}\ \prec\
  \underbrace{+\,\mathrm{SNF}}_{\mathrm{Lar}}\ \prec\
  \underbrace{+\,|\text{brackets}|}_{\mathrm{Lbr}}\ \prec\
  \underbrace{+\,\mathrm{Los}}_{\text{signed chirotope}},
\end{equation}
each layer a permutation- and $GL(5,\mathbb Z)$-invariant of $A$ computed by exact
integer linear algebra (facet normals as integer cofactor cross products; Smith
forms; brackets $|\det(\widehat a_{i_0},\dots,\widehat a_{i_5})|$). Layers
$\mathrm{Lar},\mathrm{Lbr},\mathrm{Los}$ are the computable surrogates for the
decorated Gale datum~\eqref{eq:galedatum}: $\mathrm{Lbr}$ is the
\emph{unsigned} chirotope (GKZ bracket magnitudes), $\mathrm{Los}$ the
\emph{signed} chirotope, and $\mathrm{Lar}$ the arithmetic decoration. The
signed-chirotope layer $\mathrm{Los}$ is built to be invariant under the
\emph{full} symmetry, including the global reorientation a $\det=-1$ lattice map
induces: vertices are color-ordered by their $GL$-invariant pairing signature, the
signed integer bracket is recorded on color-distinct $6$-subsets (magnitude only
on ties, where the sign is not canonical), and the resulting multiset is
symmetrized over a single global sign flip.

\begin{proposition}[Well-definedness, verified]\label{prop:keywd}
Every layer in~\eqref{eq:keylayers} is constant on $GL(5,\mathbb Z)\times
S_{n_v}$-orbits, \emph{including orientation-reversing} ($\det=-1$) maps. Hence
the key never splits a true fibre: all weight systems of one fibre receive one
key. \emph{Tested:} randomized $GL(5,\mathbb Z)$ (both orientations) $+$ vertex
relabelings on configurations across $g=1,\dots,30$, with zero changes of any
layer; the exact-arithmetic path was confirmed on scrambles with coordinates up
to $\sim\!2000$ (where floating-point determinants fail).
\end{proposition}

\subsubsection*{Completeness against the true NF: the cascade result}

A single shared-memory pass (one compute node, $64$ cores, $\approx\!26$ min)
first groups all $N$ fibres by the \emph{free} invariants already in the table
($h^{1,1},h^{1,2},h^{1,3},h^{2,2},\chi$, and lattice-point / vertex / facet
counts), then computes the geometric layers \emph{only} on the residual
collisions. The result, compared bit-for-bit against the stored normal forms:
\begin{center}\small
\begin{tabular}{l r}
\hline
quantity & value \\
\hline
fibres (distinct NF) & $88{,}588{,}676$ \\
separated by free invariants alone & $77{,}379{,}106$ \ ($87.35\%$) \\
distinct free-invariant classes & $82{,}325{,}277$ \\
needing geometry (free-collisions) & $11{,}209{,}570$ \ ($12.65\%$) \\
distinct classes of cheap key $\mathrm{(free}\oplus\mathrm{L1}\oplus\mathrm{Lpm})$ & $88{,}588{,}633$ \\
\textbf{residual collisions (over-merged fibres)} & $\mathbf{85}$ \ ($9.6\times10^{-7}$) \\
nonreflexive (sanity) & $0$ \\
\hline
\end{tabular}
\end{center}
So the cheap combined key --- the cumulative tuple through $\mathrm{Lpm}$,
i.e.\ \emph{combinatorial type $+$ pairing matrix}, requiring no Smith forms or
brackets --- is \textbf{injective on $99.99990\%$ of all fibres}. This is the
sieve of Section~\ref{sec:fibre-sieve}, stages S5--S6, made concrete and measured:
the ``near-injective signature'' S5 is in fact injective up to a $10^{-6}$
residue, and the Gale datum has replaced $\beta$ everywhere except on $85$
fibres.

\subsubsection*{The $85$-fibre hard core: a sharpening of
Theorem~\ref{thm:galerefine}}

The $85$ residual fibres fall into $42$ groups ($41$ pairs and one triple),
concentrated at low Gale dimension $g\in\{2,\dots,11\}$ (none higher). We
resolved their status exactly.

\begin{theorem}[The bounded-invariant hard core is genuine and survives the
oriented matroid]\label{thm:hardcore}
For each of the $42$ groups, the members are pairwise
$GL(5,\mathbb Z)$-\emph{inequivalent} (verified by exhaustive color-matched
unimodular search: no integer $U$ with $\det U=\pm1$ maps one vertex set onto
another), yet they are \emph{identical} on every bounded multiset invariant we
compute, namely:
\begin{enumerate}
  \item[\textnormal{(i)}] the combinatorial type $\mathrm{L1}$ and the
        vertex--facet pairing data $\mathrm{Lpm}$;
  \item[\textnormal{(ii)}] the arithmetic decoration $\mathrm{Lar}$ (Smith normal
        forms of the vertex and normal matrices);
  \item[\textnormal{(iii)}] the unsigned chirotope $\mathrm{Lbr}$ (multiset of
        bracket magnitudes); and --- most strikingly ---
  \item[\textnormal{(iv)}] the \emph{signed} chirotope $\mathrm{Los}$ (the joint
        sign-and-magnitude integer-bracket multiset, color-keyed, up to vertex
        relabeling and global reorientation).
\end{enumerate}
Consequently no layer of~\eqref{eq:keylayers} separates them: they are an
irreducible residue of the entire \emph{permutation-invariant} Gale datum.
\end{theorem}

\begin{remark}[Reading the hard core: where the program's completeness claim is
and is not literal]\label{rmk:hardcore}
Theorem~\ref{thm:galerefine} and Remark~\ref{rmk:galeprogram} state that the pair
$(\chi_B,\text{decoration})$ is a \emph{complete} $GL(5,\mathbb Z)$-invariant.
Theorem~\ref{thm:hardcore} pins down the fine print: completeness holds for the
chirotope taken \emph{up to a canonical relabeling} (the full ordered
Grassmann--Plücker sign data), but \emph{not} for any
\emph{permutation-invariant multiset summary} of it. Two
$GL(5,\mathbb Z)$-inequivalent configurations can share the same oriented matroid
\emph{as an unlabeled object} together with all elementary-divisor data; what
distinguishes them is the realization \emph{up to a canonical labeling}, and
recovering that canonical labeling is exactly the
$GL(5,\mathbb Z)$-isomorphism problem the normal form $\beta$ solves. The
$85$-fibre residue is the concrete, quantified appearance of the
``within-chirotope'' residue anticipated abstractly in
Remark~\ref{rmk:ceiling}; we find it is small ($9.6\times10^{-7}$) but
\emph{nonempty even after the full signed oriented matroid}, and present already
at $g=2$. (A methodological caution: a naive signed-bracket tally taken under the
\emph{stored NF's own coordinate order} appears to separate $37$ of the $42$
groups, but that order is not $GL$-invariant, so the apparent separation is an
artifact of comparing two different representatives, not a genuine
oriented-matroid difference --- under any honestly invariant signed-chirotope
statistic the groups do \emph{not} separate.)
\end{remark}

\begin{remark}[A complete, validated classifier with a $10^{-6}$ canonical-form
core]\label{rmk:completeclassifier}
The verification closes into a complete pipeline that uses the expensive
canonicalization only where it is provably necessary:
\[
  \text{(cheap key }\mathrm{free}\oplus\mathrm{L1}\oplus\mathrm{Lpm}\text{)}
  \ +\ \text{(exact } GL(5,\mathbb Z)\text{-isomorphism resolver on the residual
  candidates).}
\]
The exact resolver --- a finite color-restricted unimodular search, the
realization of stage S6 without invoking PALP --- splits the $42$ over-merged
classes into their $85$ constituent fibres, making the combined classifier
\emph{injective on all $88{,}588{,}676$ fibres}. Thus the empirical verdict
matches the theory exactly: the Gale datum plus one hull replaces $\beta$ for all
but a $\sim\!10^{-6}$ fraction, and that fraction is the irreducible
lattice-isomorphism core of Theorem~\ref{thm:obstruction}, here exhibited,
counted, and resolved. (Encoding note, consistent with
Remark~\ref{rmk:galecalib-read}: \texttt{nf\_vertices} is the $d\times n_v$ PALP
matrix stored column-major \texttt{int32}; the true vertex and facet counts are
\texttt{bh\_mv} and \texttt{bh\_nv}, while \texttt{vertex\_count}/%
\texttt{facet\_count} are corrupt and were not used.)
\end{remark}

%==============================================================
\section{Key references}
%==============================================================
\begin{itemize}
  \item \textbf{Kreuzer--Skarke (2000)}, ``Complete classification of
        reflexive polyhedra in four dimensions,'' \emph{Adv.\ Theor.\ Math.\
        Phys.}\ \textbf{4}, 1209--1230 (arXiv:hep-th/0002240): the
        $473\,800\,776$ enumeration and the max/min, subpolytope, and
        weight-system machinery (Sections~\ref{sec:sub}--\ref{sec:weights}).
  \item \textbf{Kreuzer--Skarke (1998)}, ``Classification of reflexive
        polyhedra in three dimensions,'' \emph{Adv.\ Theor.\ Math.\ Phys.}\
        \textbf{2}, 853--871 (arXiv:hep-th/9805190): the $4319$ enumeration
        and the weight-system / IP method in $d=3$.
  \item \textbf{Kreuzer--Skarke (2004)}, ``PALP: A Package for Analysing
        Lattice Polytopes \dots,'' \emph{Comput.\ Phys.\ Commun.}\
        \textbf{157}, 87--106 (arXiv:math/0204356): the software computing
        $GL(d,\mathbb{Z})$ normal forms used in step (5) and in
        \texttt{refpoly-5d}. The source in \texttt{PALP/} is the object of
        Section~\ref{sec:algo}: weight/CWS enumeration (\texttt{cws.c}),
        polytope construction (\texttt{Coord.c}), normal forms
        (\texttt{Polynf.c}), and subpolytope/database handling
        (\texttt{class.c}, \texttt{Subpoly.c}, \texttt{Subdb.c}); the
        \texttt{CWS} struct and size constants are in \texttt{Global.h}.
  \item \textbf{Batyrev (1994)}, ``Dual polyhedra and mirror symmetry
        \dots,'' \emph{J.\ Algebraic Geom.}\ \textbf{3}, 493--535
        (arXiv:alg-geom/9310003): reflexive polytopes, polar duality, and
        the Gorenstein-Fano / Calabi--Yau dictionary (used throughout).
  \item \textbf{Hensley (1983)}, ``Lattice vertex polytopes with interior
        lattice points,'' \emph{Pacific J.\ Math.}\ \textbf{105}, 183--191:
        the volume bound, Theorem~\ref{thm:hensley}.
  \item \textbf{Lagarias--Ziegler (1991)}, ``Bounds for lattice polytopes
        containing a fixed number of interior points in a sublattice,''
        \emph{Canad.\ J.\ Math.}\ \textbf{43}, 1022--1035: finiteness up to
        $GL(d,\mathbb{Z})$, Theorem~\ref{thm:finite}.
  \item \textbf{Sch\"oller--Skarke (2019)}, ``All Weight Systems for
        Calabi--Yau Fourfolds from Reflexivity,'' \emph{Commun.\ Math.\ Phys.}\
        (arXiv:1808.02422): the $322\,383\,760\,930$ IP /
        $185\,269\,499\,015$ reflexive single weight systems for \emph{maximal}
        reflexive $5$-polytopes, the dataset underlying
        Section~\ref{sec:fibre} (Remarks~\ref{rmk:datasetnumbers},
        \ref{rmk:fibredata}) and Theorem~\ref{thm:maxdeg5}.
  \item \textbf{Conrads (2002)}, ``Weighted projective spaces and reflexive
        simplices,'' \emph{Manuscripta Math.}\ \textbf{107}, 215--227: the
        classification of reflexive simplices by weight systems modulo an
        explicit finite equivalence --- the exact predictor for the simplex
        sub-class (Theorem~\ref{thm:dichotomy}(S),
        Section~\ref{sec:fibre-sieve}).
  \item \textbf{Gelfand--Kapranov--Zelevinsky (1994)}, \emph{Discriminants,
        Resultants, and Multidimensional Determinants}, Birkh\"auser: the
        secondary polytope / secondary fan and the $A$-discriminant
        underpinning the Gale-dual program of Section~\ref{sec:fibre-gale}
        (Remark~\ref{rmk:surrounding}).
  \item \textbf{Gr\"unbaum (2003)}, \emph{Convex Polytopes}, 2nd ed., Springer
        (GTM 221): Gale diagrams and the classification of polytopes with few
        vertices ($d{+}2$, $d{+}3$) --- the $g\le2$ strata of
        Theorem~\ref{thm:ceiling}.
  \item \textbf{Ziegler (1995)}, \emph{Lectures on Polytopes}, Springer
        (GTM 152): Gale transforms/diagrams and the Gale evaluation criterion
        for faces (Proposition~\ref{prop:wsgale},
        Definition~\ref{def:galedatum}).
  \item \textbf{Bj\"orner--Las Vergnas--Sturmfels--White--Ziegler (1999)},
        \emph{Oriented Matroids}, 2nd ed., Cambridge: chirotopes, realizability,
        and the combinatorial backbone of the refinement
        chain~\eqref{eq:refinechain} (Theorem~\ref{thm:galerefine}).
\end{itemize}

\end{document}
